\documentclass[11pt,a4paper]{article}
\usepackage{fontspec}
\setmainfont{Georgia}
\newfontfamily\georgiahead{Georgia}
\usepackage[spanish,es-tabla]{babel}
\usepackage[left=2cm, right=2cm, top=2cm, bottom=2cm, footskip=0.8cm]{geometry}
\linespread{1.08}
\usepackage{graphicx}
\usepackage{amsmath,amssymb,amsfonts}
\usepackage{xcolor}
\usepackage{fancyhdr}
\usepackage{booktabs}
\usepackage{tabularx}
\usepackage{array}
\usepackage{float}
\usepackage{multicol}
\usepackage{tcolorbox}
\tcbuselibrary{skins, breakable}
\usepackage{hyperref}
\usepackage{titlesec}
\usepackage{caption}
\usepackage{setspace}
\usepackage{microtype}
\usepackage{tocloft}

% Paleta de colores sobria y profesional
\definecolor{gsnavy}{HTML}{0D233A}
\definecolor{gsblue}{HTML}{1E3A8A}
\definecolor{riskred}{HTML}{B91C1C}
\definecolor{aluargreen}{HTML}{00843D}
\definecolor{aluargold}{HTML}{D97706}
\definecolor{gsgray}{HTML}{4A4A4A}
\definecolor{darkslate}{HTML}{374151}
\definecolor{lightgray}{HTML}{F8FAFC}
\definecolor{bordergray}{HTML}{D1D5DB}

\hypersetup{
    colorlinks=true,
    linkcolor=gsnavy,
    citecolor=gsnavy,
    urlcolor=gsblue,
    pdftitle={Aluar Aluminio Argentino - Informe de Valuación y Análisis de Riesgo},
    pdfauthor={Federico Agustín Chillón},
    pdfsubject={Valuación Fundamental por Flujo de Fondos Descontados y Modelos Estocásticos},
    pdfkeywords={Valuación, Aluar, ALUA.BA, DCF, WACC, Opciones Reales, Monte Carlo, GARCH, Cópula t-Student, Chillón},
    pdfcreator={Federico Agustín Chillón},
    pdflang={es-AR}
}

% Formato de títulos con espaciado respirado
\titleformat{\section}{\georgiahead\large\bfseries\color{black}}{\thesection}{0.8em}{}
\titleformat{\subsection}{\georgiahead\normalsize\bfseries\color{black}}{\thesubsection}{0.6em}{}
\titleformat{\subsubsection}{\georgiahead\small\bfseries\color{gsgray}}{\thesubsubsection}{0.5em}{}
\titlespacing*{\section}{0pt}{9pt plus 2pt minus 1pt}{3.5pt}
\titlespacing*{\subsection}{0pt}{7pt plus 2pt minus 1pt}{2.5pt}
\titlespacing*{\subsubsection}{0pt}{5pt plus 1pt minus 1pt}{2pt}

\newcommand{\coverheading}[1]{{\georgiahead\bfseries\color{gsnavy}\normalsize #1}\par\vspace{2pt}}

% Configuración de Captions
\captionsetup{labelfont={bf,color=gsnavy}, textfont={small,color=gsgray}, skip=4pt}

% Configuración de TOC
\renewcommand{\cftsecleader}{\cftdotfill{\cftdotsep}}
\renewcommand{\cftsecfont}{\scriptsize\bfseries\color{gsnavy}}
\renewcommand{\cftsecpagefont}{\scriptsize\bfseries\color{gsnavy}}
\renewcommand{\cftsubsecfont}{\tiny\color{gsgray}}
\renewcommand{\cftsubsecpagefont}{\tiny\color{gsgray}}
\setlength{\cftbeforesecskip}{1.5pt}
\setlength{\cftbeforesubsecskip}{0.8pt}

% Encabezados y Pie de Página
\pagestyle{fancy}
\fancyhf{}
\setlength{\headheight}{14pt}
\lhead{\small \textbf{ALUAR S.A.I.C. (ALUA.BA)}}
\rhead{\small Federico Agustín Chillón}
\lfoot{\color{gsgray}\scriptsize \textsf{Federico Agustín Chillón | federicoagustinchillon@gmail.com}}
\rfoot{\color{gsgray}\scriptsize \textsf{Página \thepage}}
\renewcommand{\headrulewidth}{0.25pt}
\renewcommand{\footrulewidth}{0.25pt}
\renewcommand{\headrule}{\hbox to\headwidth{\color{bordergray}\leaders\hrule height \headrulewidth\hfill}}
\renewcommand{\footrule}{\hbox to\headwidth{\color{bordergray}\leaders\hrule height \footrulewidth\hfill}}

\setstretch{1.10}
\setlength{\parindent}{1.2em}
\setlength{\floatsep}{8pt plus 2pt minus 2pt}
\setlength{\textfloatsep}{10pt plus 2pt minus 2pt}
\setlength{\intextsep}{8pt plus 2pt minus 2pt}
\setlength{\parskip}{2.5pt plus 1pt minus 0.5pt}
\renewcommand{\arraystretch}{1.06}

\begin{document}
\renewcommand{\topfraction}{0.95}
\renewcommand{\bottomfraction}{0.95}
\renewcommand{\floatpagefraction}{0.85}
\renewcommand{\textfraction}{0.05}

% ═════════════════════════════════════════════════════════════════════════════
% HOJA 1: PORTADA
% ═════════════════════════════════════════════

\noindent
\begin{minipage}[b]{0.65\textwidth}
    {\fontsize{22}{26}\selectfont \color{gsnavy} \textbf{Aluar Aluminio Argentino}} \\[0.15cm]
    {\large \color{gsgray} \textsf{Informe de Valuación y Análisis de Riesgo}}
\end{minipage}%
\begin{minipage}[b]{0.35\textwidth}
    \raggedleft
    {\small \color{gsgray} Análisis independiente} \\
    {\small \color{gsgray} Fecha de Valuación: 23 de julio de 2026}
\end{minipage}

\vspace{0.15cm}
\noindent\textcolor{gsnavy}{\rule{\textwidth}{1.2pt}}
\vspace{0.08cm}

\noindent
{\scriptsize \color{gsgray} Documento de análisis financiero independiente. Para el descargo de responsabilidad legal y certificaciones (Ley N.\textsuperscript{o}~26.831), consulte la página 2. Datos de mercado (LME, EMBI+, CCL) al 23 de julio de 2026. No reflejan movimientos posteriores a esa fecha.}

\vspace{0.15cm}

\noindent
\begin{minipage}[t]{0.53\textwidth}
    \coverheading{Resumen Ejecutivo}
    {\small Valuación fundamental de \textbf{Aluar Aluminio Argentino S.A.I.C. (ALUA.BA)} con dictamen de \textbf{COMPRAR} y precio objetivo fundamental base de \textbf{ARS 1.268,47} por acción (+29,1\% de retorno esperado frente a una cotización de mercado de ARS 982,50). Sumando el valor de la opción real de expansión eólica de la Etapa V del Parque Eólico Aluar (+ARS 105,74), el precio objetivo integrado llega a ARS 1.374,22 (+39,9\%). La cotización actual descuenta un crecimiento terminal implícito deprimido ($g^* = 0,69\%$), subestimando la ventaja en costos por energía propia y el desapalancamiento financiero proyectado.}

    \vspace{0.15cm}
    \coverheading{Pilares de la Tesis de Inversión}
    \begin{enumerate}
        \setlength{\itemsep}{2pt}
        \item \textbf{Estructura de Costos en Primer Cuartil Global:} Matriz energética 96\% autogenerada (78\% renovable vía Futaleufú y PEAL), con costo en efectivo C1 en USD 1.680/t (percentil 29\% mundial).
        \item \textbf{Ingresos en Moneda Dura con Costos Fijos en Pesos:} 83,8\% de las ventas se exportan a precios fijados en el LME, mientras que los costos operativos locales se licúan en términos reales ante depreciaciones del tipo de cambio.
        \item \textbf{Compresión del Diferencial Soberano y Desapalancamiento:} La caída del spread EMBI+ (441 pb) sitúa el WACC en 6,95\% ($\lambda=0,1617$), y el cese de inversiones intensivas reduce la deuda neta a $<1,0\text{x}$ EBITDA hacia 2028e.
    \end{enumerate}
\end{minipage}%
\hfill
\begin{minipage}[t]{0.45\textwidth}
    \vspace{-0.1cm}
    \begin{tcolorbox}[colback=white, colframe=gsnavy, boxrule=0.8pt, arc=3pt, left=7pt, right=7pt, top=5pt, bottom=5pt]
        \begin{center}
            {\scriptsize \color{gsgray} Dictamen del Modelo Fundamental} \\[0.08cm]
            {\fontsize{15}{18}\selectfont \sffamily \textbf{\color{aluargreen} COMPRAR}} \\[0.08cm]
            {\small \color{gsnavy} \textbf{Precio Objetivo Fundamental Base: ARS 1.268,47} (USD 0,80)} \\[0.03cm]
            {\scriptsize \color{gsgray} Cotización de Mercado: ARS 982,50 $\mid$ \textbf{Retorno: +29,1\%}} \\[0.05cm]
            \color{bordergray}\rule{0.85\linewidth}{0.3pt} \\[0.05cm]
            {\scriptsize \color{gsgray} Escenario Expansión (+ opción real de la Etapa V eólica):} \\[0.02cm]
            {\footnotesize \color{gsnavy} Precio Objetivo Integrado: ARS 1.374,22 (USD 0,84 $\mid$ +39,9\%)}
        \end{center}
        \vspace{0.08cm}
        \color{gsnavy}\rule{\linewidth}{0.4pt}
        \vspace{0.06cm}
        {\scriptsize
        \renewcommand{\arraystretch}{1.25}
        \centering
        \begin{tabular}{@{}l@{\hspace{4pt}}r@{\hspace{6pt}}l@{\hspace{4pt}}r@{}}
            WACC:                   & 6,95\%  & Ke ($\lambda=0{,}1617$): & 9,13\%  \\
            Beta (Hamada):          & 0,888   & ERP:                   & 4,18\%  \\
            Prima País ($\lambda$): & 0,71\%  & Kd post-imp.:          & 2,47\%  \\
            Estructura D/E:          & 0,486   & Crec. Terminal ($g$):  & 2,00\%  \\
            Potencial Base / Int.:   & \multicolumn{3}{r}{+29,1\% / +39,9\%}     \\
            CCL Implícito:           & \multicolumn{3}{r}{1.584,25}                       \\
        \end{tabular}
        }
        \vspace{0.06cm}
        \color{bordergray}\rule{\linewidth}{0.4pt}
        \vspace{0.04cm}
        {\scriptsize \color{gsgray} \raggedright Valuación según Dumrauf (Cap.~14), con prima de riesgo país $\lambda \times \text{EMBI}+$ y opción real por LSMC. Datos al 23-jul-2026.\par}
    \end{tcolorbox}
\end{minipage}

\vspace{0.15cm}
\noindent
\begin{tcolorbox}[colback=white, colframe=gsblue!40, boxrule=0.5pt, top=2pt, bottom=2pt, left=4pt, right=4pt, title={\small\sffamily \textbf{Divergencias Frente al Consenso de Mercado}}, coltitle=gsnavy]
    \footnotesize
    \begin{itemize}
        \setlength{\itemsep}{0pt}
        \item \textbf{Normalización del EBITDA vs. Múltiplo Deprimido:} La cotización actual descuenta el ejercicio recesivo 2025 (13,4x EV/EBITDA). La recuperación operativa proyectada para 2026e (EBITDA USD 391,2 MM) sitúa el múltiplo forward en 5,2x.
        \item \textbf{Desapalancamiento Acelerado Post-CAPEX:} Concluida la etapa intensiva del parque eólico, la generación operativa reduce la Deuda Neta a $<1,0\text{x}$ EBITDA hacia 2028e, habilitando mayor distribución de dividendos.
        \item \textbf{Cobertura Operativa y Factor $\lambda = 0,1617$:} El mercado sobreestima la prima soberana al asumir $\lambda=1,0$. La estructura con 83,8\% de ventas dolarizadas y costos operativos en moneda local fundamenta un costo de capital de 6,95\%.
    \end{itemize}
\end{tcolorbox}

\clearpage

% ═════════════════════════════════════════════════════════════════════════════
% HOJA 2: AVISO LEGAL Y DECLARACIÓN ACADÉMICA — TRABAJO FINAL UNCUYO
% ═════════════════════════════════════════════════════════════════════════════

\begin{center}
    {\Large\georgiahead\bfseries\color{gsnavy} Aviso Legal y Declaración de Responsabilidad} \\[0.2cm]
    {\small\color{gsgray} Análisis independiente elaborado por Federico Agustín Chillón}
\end{center}
\vspace{0.05cm}
\hrule height 0.8pt
\vspace{0.15cm}

\noindent
\begin{tcolorbox}[colback=white, colframe=gsnavy, boxrule=0.8pt, arc=3pt, left=8pt, right=8pt, top=6pt, bottom=6pt]
    \small
    \textbf{\color{gsnavy} 1. Naturaleza y Propósito del Documento}\\[2pt]
    {Este trabajo de valuación fundamental fue elaborado de forma personal e independiente por Federico Agustín Chillón, con fines analíticos propios sobre Aluar Aluminio Argentino S.A.I.C. (ALUA.BA). No es un informe de investigación de valores regulado por la CNV, no constituye asesoramiento financiero profesional, y el autor no posee matrícula de agente o analista registrado ante la CNV (Ley N.\textsuperscript{o}~26.831, Art. 22 y RG 622/2013).\par}

    \vspace{0.12cm}
    \textbf{\color{gsnavy} 2. Ausencia de Oferta Pública y No Asesoramiento}\\[2pt]
    {Este documento no constituye una oferta de compra o venta de valores ni una solicitud de transacciones, ni asesoramiento financiero personalizado en los términos de la Ley N.\textsuperscript{o}~26.831 y normas de la Comisión Nacional de Valores (CNV). Los inversores deben evaluar sus propios riesgos y consultar a sus asesores profesionales antes de tomar decisiones de inversión.\par}

    \vspace{0.12cm}
    \textbf{\color{gsnavy} 3. Declaración de Independencia y Conflictos de Interés}\\[2pt]
    {El autor declara lo siguiente.
    \setlength{\itemsep}{0pt}\setlength{\parsep}{0pt}\setlength{\topsep}{2pt}
    \begin{enumerate}
        \item No mantiene una posición directa en acciones de Aluar Aluminio Argentino S.A.I.C., ni opciones o derivados sobre ALUA.BA. Podría tener una exposición indirecta e inmaterial a través de fondos comunes de inversión, sobre cuya composición no tiene control ni conocimiento detallado.
        \item No recibe compensación, remuneración ni beneficio alguno vinculado a las conclusiones de este trabajo.
        \item No existe relación comercial, societaria, laboral, de asesoramiento, ni vínculo familiar con Aluar S.A.I.C., sus controlantes (Familia Madanes, CIPAL S.A., Estancias San Javier de Catamarca S.A., Angerona Group Trust), sociedades vinculadas (Hidroeléctrica Futaleufú S.A., Transpa S.A.), ni colocadores o brokers intervinientes en emisiones recientes.
        \item El presente trabajo se elabora de forma personal e independiente, sin mandato de terceros.
    \end{enumerate}\par}
\end{tcolorbox}

\vspace{0.15cm}
\noindent
{\scriptsize \color{gsgray} \raggedright Las fuentes de información, la metodología cuantitativa y la vigencia de este informe se detallan en el Anexo de Fuentes y Metodología, antes de la bibliografía.\par}

\vfill
\noindent
{\scriptsize \color{gsgray} \raggedright Ficha del Documento: Análisis independiente. Autor: Federico Agustín Chillón. Mendoza, Argentina. Fecha de cierre de datos: 23 de julio de 2026. Contacto: \texttt{federicoagustinchillon@gmail.com}.\par}

\clearpage

% ═════════════════════════════════════════════════════════════════════════════
% HOJA 3: ÍNDICE GENERAL DE SECCIONES Y CONTENIDOS
% ═════════════════════════════════════════════════════════════════════════════

\begin{center}
    {\Large\georgiahead\bfseries\color{black} Índice}
\end{center}
\vspace{-0.2cm}
\hrule height 0.8pt
\vspace{0.15cm}

\begin{multicols}{2}
\footnotesize
\makeatletter
\@starttoc{toc}
\makeatother
\end{multicols}

\clearpage

% ═════════════════════════════════════════════════════════════════════════════
% SECCIÓN 1: DESCRIPCIÓN DE LA COMPAÑÍA
% ═════════════════════════════════════════════
\section{Integración Vertical y Estructura Operativa de la Compañía}

\subsection{Perfil Operativo y Capacidad Instalada}
Aluar Aluminio Argentino S.A.I.C. es el único productor primario de aluminio integrado en Argentina y uno de los mayores complejos industriales de Sudamérica. Fundada en 1974, la compañía opera su planta de reducción electrolítica en Puerto Madryn (Provincia del Chubut) con una capacidad instalada nominal de 460.000 toneladas anuales (Figura \ref{fig:01}). Su portafolio abarca aluminio primario líquido y sólido (lingotes estándar, lingotes T, barrotes de extrusión, alambrón y placas de laminación), junto con una división de semielaborados y elaborados de mayor valor agregado (perfiles extruidos y laminados para construcción, transporte y embalaje). 

Alrededor del 83,8\% del volumen físico producido se destina a la exportación (Estados Unidos, Brasil, Europa y Japón), liquidado en dólares estadounidenses sobre cotizaciones del London Metal Exchange (LME) más primas regionales, lo que otorga a la firma una cobertura cambiaria estructural.

\subsection{Estructura Accionaria y Control Corporativo}
Según la Memoria Anual y Estados Financieros al 30-jun-2025, el capital social está integrado por 2.800 millones de acciones ordinarias escriturales de \$1 valor nominal y un voto por acción, listadas en BYMA bajo la especie ALUA.BA. Las sociedades vinculadas a la familia Madanes (CIPAL S.A., Estancias San Javier de Catamarca S.A. y Angerona Group Trust) concentran en conjunto el 45,98\% de los derechos de voto (Figura \ref{fig:02}), conformando el grupo de control frente al 54,02\% distribuido entre la ANSES-FGS (27,47\%) y el capital flotante en el mercado doméstico e internacional.

\begin{figure}[H]
\centering
\begin{minipage}[t]{0.485\textwidth}
    \centering
    \includegraphics[width=\linewidth,height=4.6cm,keepaspectratio]{../03_Modelo_y_Codigo/figuras/figura_01.png}
    \caption{\textbf{Evolución de Aluar (1974--2026)}: Capacidad instalada y etapas de expansión.}
    \label{fig:01}
\end{minipage}\hfill
\begin{minipage}[t]{0.485\textwidth}
    \centering
    \includegraphics[width=\linewidth,height=4.6cm,keepaspectratio]{../03_Modelo_y_Codigo/figuras/figura_02.png}
    \caption{\textbf{Estructura Accionaria}: Concentración del 45,98\% en el grupo de control (Familia Madanes) y 54,02\% flotante/ANSES.}
    \label{fig:02}
\end{minipage}
\end{figure}

\subsection{Matriz Energética y Posición de Costos C1}
La fundición de aluminio primario es un proceso electrointensivo donde la energía representa entre el 30\% y 40\% del costo total de producción. La ventaja estructural en costos descansa en tres ejes:
\begin{itemize}
    \item \textbf{Autonomía Energética Propia (96\% del Consumo):} La planta de Puerto Madryn es abastecida en un 54\% por la Central Hidroeléctrica Futaleufú (320 MW asignados contractualmente bajo concesión), 24\% por el Parque Eólico Aluar (PEAL, 200 MW en operación) y 18\% por ciclos combinados térmicos propios a gas natural. Sólo el 4\% restante se toma de la red nacional (SADI/CAMMESA), aislando a la compañía de aumentos en las tarifas eléctricas spot.
    \item \textbf{Posición en el Primer Cuartil de Costos Globales:} Con un costo en efectivo C1 de USD 1.680 por tonelada (percentil 29\% de la curva de costos de la industria), Aluar se ubica entre los productores más eficientes del mundo, manteniendo márgenes operativos positivos en fases contractivas del ciclo internacional.
    \item \textbf{Terminal Portuaria y Logística de Insumos:} Muelle propio de aguas profundas en Puerto Madryn para la descarga directa de alúmina importada (insumo del cual Aluar importa el 100\%, requiriendo $\sim$1,93 t de alúmina por tonelada de aluminio, indexado al \textit{Platts Alumina Index} PAX) y embarque de exportaciones.
\end{itemize}

\subsection{Certificación ASI y Mezcla de Elaborados}
Los objetivos estratégicos comprenden: (1) alcanzar el 100\% de abastecimiento renovable mediante la Etapa V del Parque Eólico Aluar, conocida públicamente como \emph{La Flecha}, de 312 MW anunciados (ampliados a 336 MW durante la construcción), (2) obtener la certificación ASI (\emph{Aluminium Stewardship Initiative}) para capturar primas de precio por aluminio de bajas emisiones en Europa, y (3) defender la cuota en productos elaborados. La participación de elaborados en el volumen total se contrajo temporalmente debido a la recesión del mercado interno, pasando de 6,0\% en 2023 a 3,3\% en 2025:

\begin{table}[H]
\centering
\caption{\textbf{Participación de la División Elaborados en el Volumen Total de Ventas} {\small\color{gsgray} (Memoria y EEFF Aluar 2023–2025)}}
\small
\begin{tabular}{lccc}
\toprule
\textbf{Métrica de Volumen} & \textbf{2023} & \textbf{2024} & \textbf{2025} \\
\midrule
Ventas Elaborados (t) & 23.104 & 17.766 & 13.498 \\
Ventas Totales (t) & 383.834 & 418.014 & 402.991 \\
\textbf{Participación Elaborados} & \textbf{6,0\%} & \textbf{4,3\%} & \textbf{3,3\%} \\
\bottomrule
\end{tabular}
\end{table}

\subsection{Diagnóstico Estratégico FODA}
\begin{itemize}
    \item \textbf{Fortalezas:} Único productor primario integrado en Argentina, 96\% de energía autogenerada, costos C1 en primer cuartil mundial, ingresos 83,8\% en moneda dura.
    \item \textbf{Oportunidades:} Primas comerciales por certificación ASI, beneficios impositivos del RIGI para la Etapa V eólica, compresión del diferencial de riesgo soberano.
    \item \textbf{Debilidades:} Dependencia de alúmina importada, capacidad nominal operando al 100\% (460 kt), endeudamiento transitorio por inversiones eólicas.
    \item \textbf{Amenazas:} Fluctuaciones cíclicas del LME, apreciación del tipo de cambio real que encarece costos laborales en USD, excedentes de producción en Asia.
\end{itemize}

\subsection{Desempeño Financiero Reciente (2020--2025)}
\begin{table}[H]
\centering
\caption{\textbf{Evolución de Indicadores Financieros Auditados 2020–2025 (USD MM)}}
\small
\begin{tabular}{lcccccc}
\toprule
\textbf{Métrica Financiera} & \textbf{2020} & \textbf{2021} & \textbf{2022} & \textbf{2023} & \textbf{2024} & \textbf{2025} \\
\midrule
Ingresos (Ventas Netas) & 823,5 & 513,5 & 654,7 & 647,8 & 915,9 & 1.092,6 \\
EBITDA & 116,5 & 110,2 & 208,8 & 103,5 & 204,7 & 163,3 \\
Margen EBITDA & 14,15\% & 21,47\% & 31,89\% & 15,97\% & 22,35\% & 14,94\% \\
Deuda Neta & 351,3 & 137,4 & 57,3 & 169,9 & 211,2 & 525,8 \\
Deuda Neta / EBITDA & 3,02x & 1,25x & 0,27x & 1,64x & 1,03x & 3,22x \\
\bottomrule
\end{tabular}
\end{table}

\noindent
{\tiny Nota. Cifras en USD convertidas al CCL de cierre de cada ejercicio. EBITDA = Resultado Operativo + D\&A (motor de valuación). Deuda Neta auditada según estados financieros consolidados.\par}

% ═════════════════════════════════════════════════════════════════════════════
% SECCIÓN 2: CONTEXTO MACROECONÓMICO Y MECÁNICA DE PROYECCIÓN
% ═════════════════════════════════════════════
\section{Entorno Macroeconómico y Modelo Operativo \texorpdfstring{($P \times Q$)}{(P x Q)}}

\subsection{Entorno Macroeconómico y Dinámica del Riesgo Soberano}
Las proyecciones incorporan el sendero de estabilización macroeconómica argentino (Figura \ref{fig:03}), con convergencia inflacionaria y reactivación gradual de la actividad industrial. La compresión del spread de riesgo soberano EMBI+ desde máximos de 2.400 pb en 2022 a 441 pb al cierre de julio de 2026 (Figura \ref{fig:04}) reduce de forma directa la tasa de descuento requerida para el capital.

\begin{figure}[H]
\centering
\begin{minipage}[t]{0.485\textwidth}
    \centering
    \includegraphics[width=\linewidth,height=4.6cm,keepaspectratio]{../03_Modelo_y_Codigo/figuras/figura_03.png}
    \caption{\textbf{Contexto Macroeconómico}: PBI real e inflación proyectada en Argentina (2020--2030E).}
    \label{fig:03}
\end{minipage}\hfill
\begin{minipage}[t]{0.485\textwidth}
    \centering
    \includegraphics[width=\linewidth,height=4.6cm,keepaspectratio]{../03_Modelo_y_Codigo/figuras/figura_04.png}
    \caption{\textbf{Compresión del Riesgo País (EMBI+)}: Descenso del spread soberano desde 2.400 pb a 441 pb.}
    \label{fig:04}
\end{minipage}
\end{figure}

\subsection{Volumen Saturado y Reversión del LME sobre los Ingresos}
La proyección de ventas descompone la dinámica en $Ingresos = P \times Q$:
\begin{enumerate}
    \item \textbf{Volumen Físico ($Q$):} La planta opera a plena capacidad (460.000 Tn nominales). Con un factor de utilización histórico del 94\%, el volumen despachado se proyecta constante en:
    \begin{equation}
        Q = 460.000\text{ Tn} \times 0,94 = \mathbf{432.400\text{ Tn/año}}
    \end{equation}
    \item \textbf{Precio Realizado ($P$):} En 2026e los ingresos alcanzan USD 1.591,4 MM impulsados por el precio spot del LME en el 1T26 ($USD\text{ }3.173,2/\text{t}$). Para el período 2027E--2030E, el precio realizado converge linealmente a su nivel de equilibrio de largo plazo:
    \begin{equation}
        P_{2030} = LME_{\text{rev}} + Premium_{\text{elab}} = 2.587,1 + 645,7 = \mathbf{USD\text{ }3.232,8/\text{t}}
    \end{equation}
    \item \textbf{Calibración del Margen EBITDA ($k$):} Se calibra el parámetro de eficiencia operativa $k$ sobre los datos reales del 1T26 ($Ventas = USD\text{ }416\text{ MM}$, $EBITDA = USD\text{ }106\text{ MM}$, $C1 = USD\text{ }1.680/\text{t}$):
    \begin{equation}
        k = \frac{Margen_{EBITDA, Q1}}{\left(\frac{P_{Q1} - C1}{P_{Q1}}\right)} = \frac{106 / 416}{\left(\frac{3.848,3 - 1.680}{3.848,3}\right)} = \mathbf{0,4521}
    \end{equation}
    El margen EBITDA de cada período se deduce como $Margen_t = k \times \left(\frac{P_t - 1.680}{P_t}\right)$.
    \item \textbf{Variación de Capital de Trabajo ($\Delta$NWC) en 2026e:} Se proyecta según la mediana histórica del capital de trabajo operativo sobre ventas ($NWC/\text{Ventas} = 49,42\%$):
    \begin{equation}
        \Delta NWC_{2026E} = (Ventas_{2026E} - Ventas_{2025}) \times 49,42\% = (1.591,4 - 1.092,6) \times 0,4942 = \mathbf{USD\text{ }246,5\text{ MM}}
    \end{equation}
    Esta absorción inicial de capital de trabajo obedece a la recomposición de cuentas por cobrar ($DSO \approx 45\text{ días}$) y la reposición de existencias de alúmina a precios internacionales más altos tras el shock de facturación por el LME (con volumen físico constante en 432.400 t). Luego libera $+USD\text{ }23,9\text{ MM/año}$ de forma recurrente en el resto del horizonte proyectado, a medida que el precio converge a su media histórica.
\end{enumerate}

% ═════════════════════════════════════════════════════════════════════════════
% SECCIÓN 3: DINÁMICA DE LA INDUSTRIA GLOBAL DEL ALUMINIO
% ═════════════════════════════════════════════
\section{Dinámica de la Industria Global del Aluminio y Cointegración con el LME}

\subsection{Dinámica Cíclica del Aluminio, Proceso AR(1) Soberano y Sensibilidad al DXY}
La correlación negativa entre la cotización del aluminio y el índice DXY (Figura \ref{fig:06}) expone la sensibilidad cambiaria global. Aluar cuenta con una cobertura operativa natural: en fases de apreciación global del dólar acompañadas por debilidad del metal, las correcciones del tipo de cambio real reducen los costos fijos en pesos (salarios y servicios contratados localmente), amortiguando la presión sobre los márgenes operativos. Para el diferencial soberano, se proyecta una convergencia estocástica autorregresiva mensual $\text{AR}(1)$ con $\hat{\phi} = 0,7582$ (Figura \ref{fig:05}).

\begin{figure}[H]
\centering
\begin{minipage}[t]{0.485\textwidth}
    \centering
    \includegraphics[width=\linewidth,height=4.6cm,keepaspectratio]{../03_Modelo_y_Codigo/figuras/figura_05.png}
    \caption{\textbf{Trayectoria del Riesgo País}: Proyección estocástica AR(1) frente al supuesto corporativo base (441 pb).}
    \label{fig:05}
\end{minipage}\hfill
\begin{minipage}[t]{0.485\textwidth}
    \centering
    \includegraphics[width=\linewidth,height=4.6cm,keepaspectratio]{../03_Modelo_y_Codigo/figuras/figura_06.png}
    \caption{\textbf{Aluminio LME vs. DXY}: Relación inversa entre el metal y el índice global del dólar.}
    \label{fig:06}
\end{minipage}
\end{figure}

\subsection{Test de Cointegración de Engle-Granger entre ALUA y el LME}
Se sometió la serie diaria de ALUA.BA (expresada en USD vía CCL, $N=2.355$ ruedas, período 2016--2026) y el futuro del LME a la prueba de cointegración de Engle-Granger. El estadístico $t = -2,29$ ($p = 0,38$, frente a un valor crítico de $-3,34$ al 5\%) no rechaza la hipótesis nula de no cointegración estricta en niveles. A pesar de no verificarse cointegración lineal en niveles, la elasticidad contemporánea se sitúa en 0,66 ($R^2=0,12$), confirmando que Aluar responde conjuntamente al ciclo internacional del metal y a la dinámica macroeconómica local.

\subsection{Estructura de la Oferta Global y Autonomía Energética Comparada}
La producción primaria mundial está concentrada en China (56\% del total, Figura \ref{fig:07}), donde la mayor parte de las fundiciones operan a carbón. Frente a los principales productores internacionales, Aluar cuenta con energía propia: el 96\% de su consumo es autogenerado con un 78\% de fuentes renovables (Figura \ref{fig:08}), lo que reduce sus costos unitarios.

\begin{figure}[H]
\centering
\begin{minipage}[t]{0.485\textwidth}
    \centering
    \includegraphics[width=\linewidth,height=4.6cm,keepaspectratio]{../03_Modelo_y_Codigo/figuras/figura_07.png}
    \caption{\textbf{Producción Mundial de Aluminio}: Participación de los principales países productores.}
    \label{fig:07}
\end{minipage}\hfill
\begin{minipage}[t]{0.485\textwidth}
    \centering
    \includegraphics[width=\linewidth,height=4.6cm,keepaspectratio]{../03_Modelo_y_Codigo/figuras/figura_08.png}
    \caption{\textbf{Capacidad y Autogeneración}: Comparativa de escala y matriz energética propia frente a pares.}
    \label{fig:08}
\end{minipage}
\end{figure}

\subsection{Posición Competitiva Doméstica y Curva Global de Costos en Efectivo C1}
En el mercado argentino, Aluar abastece el 60\% del consumo (Figura \ref{fig:09}) con precios fijados por paridad de importación. A nivel internacional, su costo en efectivo C1 de USD 1.680 por tonelada la ubica en el percentil 29\% de la curva global (Figura \ref{fig:10}), lo que sostiene márgenes positivos aun en bajas del ciclo de commodities. El principal riesgo de costos externos proviene de la volatilidad del PAX y fletes marítimos sobre la alúmina importada.

\begin{figure}[H]
\centering
\begin{minipage}[t]{0.485\textwidth}
    \centering
    \includegraphics[width=\linewidth,height=4.6cm,keepaspectratio]{../03_Modelo_y_Codigo/figuras/figura_09.png}
    \caption{\textbf{Participación de Mercado Local}: Cuota del 60\% de Aluar frente al 40\% importado.}
    \label{fig:09}
\end{minipage}\hfill
\begin{minipage}[t]{0.485\textwidth}
    \centering
    \includegraphics[width=\linewidth,height=4.6cm,keepaspectratio]{../03_Modelo_y_Codigo/figuras/figura_10.png}
    \caption{\textbf{Matriz Renovable y Curva C1}: Curva global de costos y posición de Aluar.}
    \label{fig:10}
\end{minipage}
\end{figure}

% ═════════════════════════════════════════════════════════════════════════════
% SECCIÓN 4: GOBERNANZA CORPORATIVA Y ESG
% ═════════════════════════════════════════════
\section{Gobernanza Corporativa, Marco Regulatorio y Factores ESG}

\subsection{Gobierno Corporativo y Transparencia (Ley 26.831 / CNV)}
Aluar cotiza en el régimen de oferta pública de la CNV (Resolución General 797/2019). El Directorio está encabezado por Javier Madanes Quintanilla. El Comité de Auditoría cuenta con mayoría de directores independientes según el Art. 109 de la Ley 26.831, dictaminando sobre el sistema de control interno y transacciones con sociedades vinculadas como Hidroeléctrica Futaleufú S.A. (60,2\%) y Transpa S.A. (20,4\%), celebradas a precios de mercado.

\begin{table}[H]
\centering
\caption{\textbf{Matriz de Evaluación de Factores Ambientales, Sociales y de Gobierno (ESG)}}
\scriptsize
\begin{tabularx}{\textwidth}{p{2.5cm}Xp{4.2cm}}
\toprule
\textbf{Dimensión} & \textbf{Indicadores Observables} & \textbf{Evaluación} \\
\midrule
\textbf{Ambiental (E)} & 78\% renovable propia. Huella $<$\,4 t CO\textsubscript{2}/t Al. Certificación ASI en proceso. & \textbf{Favorable.} Ventaja competitiva frente al mecanismo de ajuste en frontera por carbono (CBAM). \\
\textbf{Social (S)} & Principal empleador privado en Chubut ($>$\,6.200 puestos directos e indirectos). TRIR $<$\,1,0. & \textbf{Adecuada.} Estándares de seguridad industrial y baja tasa de siniestralidad. \\
\textbf{Gobernanza (G)} & Concentración del 45,98\% en el grupo de control. Auditoría externa por PwC y comités independientes. & \textbf{Estable.} Continuidad estratégica con flotante del 54,02\%. \\
\bottomrule
\end{tabularx}
\end{table}

\subsection{Impacto del Ajuste de Carbono en Frontera (CBAM de la Unión Europea)}
La reducida huella de carbono de Aluar ($<$\,4 t CO\textsubscript{2}/t Al) frente al promedio de la industria ($\sim$\,11 t CO\textsubscript{2}/t) le otorga una ventaja arancelaria neta estimada en USD 80--120 por tonelada bajo el esquema CBAM de la Unión Europea frente a productores alimentados con fuentes fósiles.

% ═════════════════════════════════════════════════════════════════════════
% SECCIÓN 5: VALUACIÓN RELATIVA Y DESAPALANCAMIENTO
% ═════════════════════════════════════════════
\section{Valuación Relativa y Perfil de Desapalancamiento}

\subsection{Comparación de Múltiplos Históricos y Proyectados}
Aluar cotiza a un múltiplo EV/EBITDA de los últimos doce meses de 13,4x sobre el resultado deprimido de 2025 (USD 163,3 MM, Figura \ref{fig:11}). Al normalizar la proyección sobre el EBITDA estimado de 2026e (USD 391,2 MM), el múltiplo implícito desciende a 5,2x EV/EBITDA proyectado, situando a la firma con un descuento del 20\%--25\% frente a comparables internacionales como Alcoa (8,5x) y Norsk Hydro (7,7x).

\subsection{Trayectoria de Desapalancamiento Financiero}
El desembolso en bienes de uso durante 2025 (USD 243,9 MM) elevó la Deuda Neta a USD 525,8 MM (3,22x EBITDA). Con la moderación de inversiones hacia niveles de mantenimiento ($\sim\text{USD } 90\text{--}103\text{ MM/año}$), el flujo libre de caja proyectado permite recomprimir la Deuda Neta / EBITDA a $<1,0\text{x}$ hacia 2028e (Figura \ref{fig:12}).

\begin{figure}[H]
\centering
\begin{minipage}[t]{0.485\textwidth}
    \centering
    \includegraphics[width=\linewidth,height=4.6cm,keepaspectratio]{../03_Modelo_y_Codigo/figuras/figura_11.png}
    \caption{\textbf{EV/EBITDA vs. Margen}: Ubicación de Aluar frente a pares internacionales (Alcoa, Hydro, Chalco).}
    \label{fig:11}
\end{minipage}\hfill
\begin{minipage}[t]{0.485\textwidth}
    \centering
    \includegraphics[width=\linewidth,height=4.6cm,keepaspectratio]{../03_Modelo_y_Codigo/figuras/figura_12.png}
    \caption{\textbf{EBITDA Histórico y Proyectado}: Recuperación operativa y desapalancamiento tras el pico de CAPEX.}
    \label{fig:12}
\end{minipage}
\end{figure}

\begin{table}[H]
\centering
\caption{\textbf{Múltiplos Financieros de Empresas Comparables del Sector (julio-2026)}}
\small
\begin{tabular}{llccc}
\toprule
\textbf{Compañía} & \textbf{Ticker} & \textbf{EV/EBITDA (x)} & \textbf{Margen EBITDA} & \textbf{Deuda Neta/EBITDA} \\
\midrule
Aluar S.A.I.C. (UDM / Proy.) & ALUA.BA & 13,4x / 5,2x & 14,9\% & 3,20x \\
Alcoa Corp. & AA & 8,5x & 13,1\% & 0,51x \\
Norsk Hydro & NHY.OL & 7,7x & 11,9\% & 0,39x \\
Chalco & 2600.HK & 5,2x & 15,5\% & 0,61x \\
\midrule
\textbf{Mediana de Pares} & & \textbf{7,7x} & \textbf{13,1\%} & \textbf{0,51x} \\
\bottomrule
\end{tabular}
\end{table}

% ═════════════════════════════════════════════════════════════════════════
% SECCIÓN 6: MODELO WACC Y FACTOR LAMBDA
% ═════════════════════════════════════════════
\section{Costo Promedio Ponderado del Capital (WACC) y Factor \texorpdfstring{$\lambda$}{lambda}}

\subsection{Formulación y Estimación del WACC}
El Costo Promedio Ponderado del Capital se calcula conforme a la formulación clásica:
\begin{equation}
    WACC = K_e \left(\frac{E}{V}\right) + K_d (1 - t) \left(\frac{D}{V}\right)
\end{equation}
Reemplazando los parámetros calibrados ($E/V = 67,29\%$, $D/V = 32,71\%$, $K_d = 3,80\%$, $t = 35\%$ y $K_e = 9,13\%$), la descomposición por fuentes de financiamiento sitúa el costo de capital de régimen en 6,95\% (Figura \ref{fig:14}):
\begin{equation}
    WACC = 9,13\% \times 0,6729 + 3,80\% \times (1 - 0,35) \times 0,3271 = 6,26\% + 0,81\% = \mathbf{6,95\%}
\end{equation}
El costo pre-tax de la deuda ($K_d = 3,80\%$ en USD) refleja la estructura financiera de Aluar: Obligaciones Negociables Clase 6 y 7 dólar linked / hard dollar en el mercado local (tasas del 3,00\%–4,50\%) y líneas de prefinanciación de exportaciones garantizadas con contratos \textit{off-take}.

\subsection{Deducción del Beta en Cuatro Pasos (OLS \texorpdfstring{$\to$}{->} Blume \texorpdfstring{$\to$}{->} Hamada)}
La secuencia metodológica de estimación del riesgo sistemático (Figura \ref{fig:13}) desagrega:
\begin{enumerate}
    \item \textbf{Beta OLS Histórico:} Regresión lineal diaria de 10 años ($N=2.376$ ruedas) de ALUA (en USD CCL) contra el índice S\&P 500:
    \begin{equation}
        r_{ALUA, t} = \alpha + \beta_{OLS} r_{SP500, t} + \epsilon_t \implies \beta_{OLS} = 0,8420 \quad (SE = 0,0563, \; R^2 = 8,60\%)
    \end{equation}
    \item \textbf{Ajuste de Blume (1975) -- Corrección por Reversión a la Media Sistemática:} Corrige el error de muestreo y la tendencia empírica del riesgo sistemático a revertir hacia la media del mercado ($\beta = 1,0$) a largo plazo:
    \begin{equation}
        \beta_{Blume} = \frac{2}{3} \beta_{OLS} + \frac{1}{3} (1,0) = \frac{2}{3} (0,8420) + \frac{1}{3} = \mathbf{0,8947}
    \end{equation}
    \item \textbf{Desapalancamiento por Hamada (1972) -- Aislamiento del Riesgo Operativo Puro:} Aísla el riesgo operativo puro del negocio ($\beta_U = 0,6745$) quitándole el apalancamiento financiero histórico transitorio ($D/E_{hist} = 0,5036$):
    \begin{equation}
        \beta_U = \frac{\beta_{Blume}}{1 + (1 - t) \left(\frac{D}{E}\right)_{hist}} = \frac{0,8947}{1 + (1 - 0,35) \times 0,5036} = \mathbf{0,6745}
    \end{equation}
    \item \textbf{Reapalancamiento al $D/E$ Objetivo -- Incorporación de la Estructura de Capital de Régimen:} Incorpora la estructura de capital de régimen ($D/E_{obj} = 0,4860$) para obtener el beta apalancado sin arrastrar distorsiones de deuda:
    \begin{equation}
        \beta_L = \beta_U \left[ 1 + (1 - t) \left(\frac{D}{E}\right)_{obj} \right] = 0,6745 \times [1 + (1 - 0,35) \times 0,4860] = \mathbf{0,8876}
    \end{equation}
\end{enumerate}

\begin{figure}[H]
\centering
\begin{minipage}[t]{0.485\textwidth}
    \centering
    \includegraphics[width=\linewidth,height=4.6cm,keepaspectratio]{../03_Modelo_y_Codigo/figuras/figura_13.png}
    \caption{\textbf{Secuencia del Beta}: OLS (0,842) $\rightarrow$ Blume (0,895) $\rightarrow$ Desapalancado (0,675) $\rightarrow$ Hamada ($\beta_L = 0,8876$).}
    \label{fig:13}
\end{minipage}\hfill
\begin{minipage}[t]{0.485\textwidth}
    \centering
    \includegraphics[width=\linewidth,height=4.6cm,keepaspectratio]{../03_Modelo_y_Codigo/figuras/figura_14.png}
    \caption{\textbf{Composición del WACC}: Costo del equity $K_e = 9,13\%$, deuda post-tax $K_d = 2,47\%$, resultando en \textbf{6,95\%}.}
    \label{fig:14}
\end{minipage}
\end{figure}

\subsection{Factor de Exposición Soberana \texorpdfstring{$\lambda = 0,1617$}{lambda = 0,1617}}
El costo del capital propio incorpora el ajuste por riesgo país bajo el modelo de Damodaran para economías emergentes, referenciado por Dumrauf (Cap. 14):
\begin{equation}
    K_e = R_f + \beta_L \times ERP + \lambda \times EMBI+ = 4,70\% + 0,8876 \times 4,18\% + 0,1617 \times 4,41\% = \mathbf{9,13\%}
\end{equation}
En lugar de adicionar el spread soberano completo ($EMBI+$), se calibra $\lambda = 0,1617$ por la cobertura natural de Aluar: el 83,8\% de sus ventas son exportaciones liquidadas en dólares a precios LME y sus costos fijos operativos están fijados en pesos locales (salarios e insumos nacionales), por lo que las depreciaciones cambiarias licúan costos en moneda dura y amortiguan los shocks de riesgo país.

\begin{table}[H]
\centering
\caption{\textbf{Sensibilidad del Costo de Capital y Precio Objetivo ante Variaciones en $\lambda$}}
\small
\begin{tabularx}{\textwidth}{Xcccc}
\toprule
\textbf{Parámetro de Riesgo País} & $\mathbf{\lambda = 0,1617\text{ (Caso Base)}}$ & $\mathbf{\lambda = 0,50}$ & $\mathbf{\lambda = 0,75}$ & $\mathbf{\lambda = 1,00\text{ (Sin Cobertura)}}$ \\
\midrule
Prima por Riesgo País ($\lambda \times EMBI+$) & 0,71\% & 2,21\% & 3,31\% & 4,41\% \\
Costo del Equity ($K_e$) & 9,13\% & 10,62\% & 11,72\% & 12,82\% \\
\textbf{WACC Resultante (USD)} & \textbf{6,95\%} & \textbf{7,95\%} & \textbf{8,70\%} & \textbf{9,44\%} \\
\textbf{Precio Objetivo Base (ARS)} & \textbf{1.268,47} & \textbf{1.020,55} & \textbf{884,78} & \textbf{775,86} \\
Retorno Esperado vs. Cotización Actual & +29,1\% & +3,9\% & -9,9\% & -21,0\% \\
\bottomrule
\end{tabularx}
\end{table}

% ═════════════════════════════════════════════════════════════════════════
% SECCIÓN 7: ANÁLISIS DE CREACIÓN DE VALOR (EVA)
% ═════════════════════════════════════════════
\section{Creación de Valor Económico (EVA) y Disciplina Terminal}

\subsection{Evolución del EVA y Efecto del Capital Invertido}
El Valor Económico Agregado ($EVA_t = (ROIC_t - WACC) \times NOA_{t-1}$) mide la generación de rentabilidad económica sobre los activos operativos netos. En 2025 el EVA se situó en $-$USD 60,3 MM debido a la incorporación de USD 243,9 MM en activos eólicos al capital invertido (NOA $= \text{USD } 1.704,3\text{ MM}$) en forma previa a su plena entrada en régimen operativo.

\begin{table}[H]
\centering
\caption{\textbf{Evolución Histórica del EVA (2020--2025) y Cierre de Régimen Terminal (USD MM)}}
\scriptsize
\begin{tabularx}{\textwidth}{Xcccccc|c}
\toprule
\textbf{Métrica de Valor} & \textbf{2020} & \textbf{2021} & \textbf{2022} & \textbf{2023} & \textbf{2024} & \textbf{2025} & \textbf{Terminal (2030E)} \\
\midrule
NOPAT & 31,3 & 39,9 & 105,4 & 40,3 & 99,2 & 60,0 & 135,5 \\
Capital Invertido (NOA) & 838,4 & 531,3 & 590,5 & 771,6 & 1.248,9 & 1.704,3 & 1.919,3 \\
ROIC (NOPAT / NOA) & 3,73\% & 7,50\% & 17,85\% & 5,22\% & 7,94\% & 3,52\% & \textbf{6,95\%} \\
WACC & 6,95\% & 6,95\% & 6,95\% & 6,95\% & 6,95\% & 6,95\% & \textbf{6,95\%} \\
Spread (ROIC - WACC) & -3,33\% & +0,44\% & +10,79\% & -1,84\% & +0,88\% & -3,54\% & \textbf{0,00\%} \\
\textbf{EVA (USD MM)} & \textbf{-27,9} & \textbf{2,3} & \textbf{63,7} & \textbf{-14,2} & \textbf{11,0} & \textbf{-60,3} & \textbf{0,0} \\
\bottomrule
\end{tabularx}
\end{table}

\subsection{Caso Base y Ajuste por Reinversión del Flujo Terminal}
\begin{itemize}
    \item \textbf{Caso Base Oficial (Dumrauf, Cap. 14):} En estado estacionario a plena capacidad (460 kt), $CAPEX = D\&A$ y $\Delta NWC = 0$, por lo que $FCFF_{\text{terminal}} = NOPAT_{2030} = \text{USD } 135,5\text{ MM}$. El valor terminal resulta $TV = \frac{135,5 \times 1,02}{0,0695 - 0,02} = \mathbf{USD\text{ }2.791,4\text{ MM}}$, que descontado a valor presente y sumado al flujo explícito determina el precio objetivo de \textbf{ARS 1.268,47}.
    \item \textbf{Ajuste por Tasa de Reinversión (Damodaran, $ROIC \to WACC$):} Con una reinversión implícita $RR = g / WACC = 2,0\% / 6,95\% = 28,78\%$, el flujo terminal ajustado resulta en $NOPAT \times (1 - RR) = \text{USD } 96,5\text{ MM}$, determinando una valuación de ARS 921,10 por acción.
\end{itemize}

% ═════════════════════════════════════════════════════════════════════════
% SECCIÓN 8: PUENTE DE VALUACIÓN
% ═════════════════════════════════════════════
\section{Puente de Valuación}

\subsection{Composición del Enterprise Value y Valor Patrimonial}
El Enterprise Value fundamental asciende a USD 2.698 MM (Figura \ref{fig:15}). Los flujos de fondos explícitos (2026e--2030e) aportan USD 564 MM (20,9\% del EV), mientras que el Valor Terminal representa USD 2.134 MM (79,1\% del EV), reflejando la vida útil indefinida de los activos de fundición. Deduciendo la Deuda Neta ajustada de USD 456,0 MM, se obtiene un Equity Value de USD 2.242 MM, equivalente a ARS 1.268,47 por acción al tipo de cambio CCL de 1.584,25.

\subsection{Evolución del Ciclo de Conversión de Efectivo (CCC)}
El Ciclo de Conversión de Efectivo (Figura \ref{fig:16}) se comprime desde 79 días en 2023 a 64 días proyectados, producto de la normalización en días de inventario (DIO) y cobranza (DSO), favoreciendo la liquidez para el servicio de pasivos.

\begin{figure}[H]
\centering
\begin{minipage}[t]{0.485\textwidth}
    \centering
    \includegraphics[width=\linewidth,height=4.6cm,keepaspectratio]{../03_Modelo_y_Codigo/figuras/figura_15.png}
    \caption{\textbf{Puente de Valuación}: De Enterprise Value (USD 2.698 MM) a Equity Value (USD 2.242 MM) y Precio Objetivo.}
    \label{fig:15}
\end{minipage}\hfill
\begin{minipage}[t]{0.485\textwidth}
    \centering
    \includegraphics[width=\linewidth,height=4.6cm,keepaspectratio]{../03_Modelo_y_Codigo/figuras/figura_16.png}
    \caption{\textbf{Capital de Trabajo Operativo}: Trayectoria de DIO, DSO, DPO y Ciclo de Conversión de Efectivo (CCC).}
    \label{fig:16}
\end{minipage}
\end{figure}

\noindent
{\tiny Nota de conciliación de deuda. La Deuda Neta contable a junio-2025 es de USD 525,8 MM (Deuda USD 594,9 MM menos Caja USD 69,2 MM). En el puente de valuación se deduce la Deuda Neta ajustada de USD 456,0 MM al inicio de las proyecciones 2026E, neta de créditos no operativos y amortizaciones del período.\par}

% ═════════════════════════════════════════════════════════════════════════
% SECCIÓN 9: ANÁLISIS ESTRATÉGICO DUPONT Y FRAMEWORKS
% ═════════════════════════════════════════════
\section{Análisis Estratégico DuPont y Fuerzas Competitivas}

\subsection{Descomposición DuPont de Cinco Factores (2020--2025)}
La descomposición DuPont extendida permite explicar la variabilidad del ROE histórico:
\begin{equation}
    ROE = \underbrace{\left(\frac{\text{RN}}{\text{EBT}}\right)}_{\text{Carga Impositiva}} \times \underbrace{\left(\frac{\text{EBT}}{\text{EBIT}}\right)}_{\text{Carga Financiera}} \times \underbrace{\left(\frac{\text{EBIT}}{\text{Ventas}}\right)}_{\text{Margen EBIT}} \times \underbrace{\left(\frac{\text{Ventas}}{\text{Activos}}\right)}_{\text{Rotación Activos}} \times \underbrace{\left(\frac{\text{Activos}}{\text{PN}}\right)}_{\text{Apalancamiento}}
\end{equation}
En 2025, el factor impositivo se redujo a 15,67\% (alícuota efectiva del 84,3\% por el impacto del ajuste por inflación impositivo NIC 29), situando el ROE en 0,83\% aun cuando el margen operativo EBIT se ubicó en 8,44\%.

\begin{table}[H]
\centering
\caption{\textbf{Descomposición DuPont de 3 Factores (2020--2025)}}
\small
\begin{tabularx}{\textwidth}{Xcccccc}
\toprule
\textbf{Componente DuPont} & \textbf{2020} & \textbf{2021} & \textbf{2022} & \textbf{2023} & \textbf{2024} & \textbf{2025} \\
\midrule
Margen Neto ($RN / Ventas$) & -5,34\% & 5,48\% & 17,65\% & 21,46\% & 9,80\% & 0,84\% \\
Rotación Activos ($V / Activos$) & 0,78x & 0,74x & 0,79x & 0,71x & 0,63x & 0,55x \\
Apalancamiento ($Activos / PN$) & 2,27x & 2,00x & 1,82x & 1,68x & 1,74x & 1,78x \\
\midrule
\textbf{ROE (Producto)} & \textbf{-9,49\%} & \textbf{8,04\%} & \textbf{25,31\%} & \textbf{25,62\%} & \textbf{10,68\%} & \textbf{0,83\%} \\
\bottomrule
\end{tabularx}
\end{table}

\subsection{Las Cinco Fuerzas de Porter}
\begin{itemize}
    \item \textbf{Amenaza de Nuevos Entrantes (Muy Baja):} Elevada inversión de capital ($>$\,USD 2.500 MM para un complejo nuevo) y acceso restringido a fuentes hidroeléctricas dedicadas.
    \item \textbf{Poder de Negociación de Proveedores (Bajo a Medio):} Concentración en alúmina importada mitigada por contratos de suministro diversificados, insumo energético 96\% autogenerado.
    \item \textbf{Poder de Negociación de Clientes (Bajo):} Exportaciones colocadas a precios internacionales tomadores LME, cuota del 60\% en el mercado doméstico.
    \item \textbf{Amenaza de Sustitutos (Baja):} Propiedades mecánicas y reciclabilidad del aluminio difíciles de reemplazar en extrusión estructural y transporte.
    \item \textbf{Rivalidad Competitiva (Media):} Competencia internacional por costos, protección operativa en el primer cuartil C1.
\end{itemize}

\subsection{Impacto del Régimen RIGI}
El Régimen de Incentivos para Grandes Inversiones (RIGI) reduce la tasa corporativa del impuesto a las ganancias al 25\% para nuevos desarrollos eólicos y elimina aranceles de importación en bienes de capital, optimizando el rendimiento de expansiones futuras.

% ═════════════════════════════════════════════════════════════════════════
% SECCIÓN 10: ANÁLISIS FINANCIERO HISTÓRICO Y RATIOS
% ═════════════════════════════════════════════
\section{Análisis Financiero Histórico, Ratios Operativos y Cobertura}

\subsection{Desempeño Operativo, Solvencia y Cobertura Financiera}
La serie financiera auditada (2020--2025) comprende cuatro ejes operativos. 1) Rentabilidad (ROIC y Márgenes): margen EBITDA promedio de 20,1\% (pico de 31,9\% en 2022) y ROIC de 3,52\% en 2025 por la asignación de USD 243,9 MM al Parque Eólico (Etapa V), convergiendo al 6,95\% terminal. 2) Solvencia: Deuda Neta / EBITDA de 3,22x en el cierre del plan de inversión, con cobertura de intereses de 3,53x ($EBIT/\text{Intereses}$) y Deuda Total / PN de 0,54x. 3) Liquidez: razón corriente de 2,44x (prueba ácida de 0,55x) para respaldar compras internacionales de alúmina. 4) Eficiencia: ciclo de conversión de efectivo en 70 días (DIO = 102, DSO = 34, DPO = 66).

\begin{table}[H]
\centering
\caption{\textbf{Estados Financieros y Ratios Operativos Históricos de Aluar (USD MM)}}
\scriptsize
\begin{tabularx}{\textwidth}{Xcccccc}
\toprule
\textbf{Métrica / Ejercicio} & \textbf{2020} & \textbf{2021} & \textbf{2022} & \textbf{2023} & \textbf{2024} & \textbf{2025} \\
\midrule
Ventas Netas & 823,50 & 513,50 & 654,72 & 647,75 & 915,85 & 1.092,60 \\
EBITDA & 116,50 & 110,23 & 208,81 & 103,47 & 204,70 & 163,25 \\
Margen EBITDA & 14,15\% & 21,47\% & 31,89\% & 15,97\% & 22,35\% & 14,94\% \\
EBIT & 48,10 & 61,33 & 162,14 & 61,97 & 152,57 & 92,27 \\
Margen EBIT & 5,84\% & 11,94\% & 24,76\% & 9,57\% & 16,66\% & 8,44\% \\
Resultado Neto & -43,94 & 28,13 & 115,53 & 139,01 & 89,74 & 9,17 \\
Margen Neto & -5,34\% & 5,48\% & 17,65\% & 21,46\% & 9,80\% & 0,84\% \\
CAPEX & -103,01 & -9,43 & -11,11 & -65,89 & -25,59 & -243,93 \\
Deuda Neta & 351,30 & 137,44 & 57,33 & 169,88 & 211,23 & 525,76 \\
\midrule
\textbf{ROIC} & \textbf{3,73\%} & \textbf{7,50\%} & \textbf{17,85\%} & \textbf{5,22\%} & \textbf{7,94\%} & \textbf{3,52\%} \\
ROE & -9,49\% & 8,04\% & 25,31\% & 25,62\% & 10,68\% & 0,83\% \\
Liquidez Corriente & 1,89x & 3,78x & 2,92x & 2,92x & 4,93x & 2,44x \\
Deuda Neta / EBITDA & 3,02x & 1,25x & 0,27x & 1,64x & 1,03x & 3,22x \\
Cobertura Intereses ($EBIT/Int.$) & 1,81x & 4,12x & 28,35x & 3,87x & 11,63x & 3,53x \\
\bottomrule
\end{tabularx}
\end{table}

% ═════════════════════════════════════════════════════════════════════════
% SECCIÓN 11: VALUACIÓN FUNDAMENTAL (DCF)
% ═════════════════════════════════════════════
\section{Valuación por Flujo de Fondos Descontados (DCF)}

\subsection{Estructura del Modelo DCF y Valor de la Opción Real de la Etapa V Eólica}
El valor patrimonial se determina mediante el descuento de flujos de fondos libres de la firma (FCFF) conforme a la metodología de Dumrauf (Cap. 14):
\begin{equation}
    EV = \sum_{t=0}^{4} \frac{FCFF_{2026+t}}{(1 + WACC)^t} + \frac{TV}{(1 + WACC)^4} = 564 + 2.134 = \mathbf{USD\text{ }2.698\text{ MM}}
\end{equation}
\begin{itemize}
    \item \textbf{Precio Objetivo DCF Caso Base:} \textbf{ARS 1.268,47} por acción (USD 0,80, +29,1\% de retorno esperado).
    \item \textbf{Opción Real sobre la Etapa V Eólica (Longstaff-Schwartz):} Cuantifica el valor de poder diferir o ejecutar la ampliación de 312 MW anunciados originalmente (312--336 MW según etapa de construcción) según evolucionen las cotizaciones del metal y el costo eléctrico, modelada mediante mínimos cuadrados (LSMC) con polinomios de Laguerre ($I_0 = \text{USD } 400\text{ MM}$, inversión estimada de la Etapa V, tenor de 5 años, $\sigma = 19,97\%$ del proceso GBM del LME calibrado sobre la serie histórica diaria de cotizaciones del metal (2016--2026), $R_f = 4,70\%$, ahorro de USD 25--30/MWh):
    \begin{equation}
        \hat{C}(S_t) = \mathbb{E}^{\mathbb{Q}}\left[ e^{-r \Delta t} V_{t+1} \mid S_t \right] = \sum_{j=0}^{2} a_j L_j(S_t)
    \end{equation}
    El subyacente $S_0$ (valor presente del ahorro incremental de energía) se modela como la perpetuidad del ahorro a la tasa WACC ($312\text{ MW} \times 45\% \times 8.760\text{ h} \times \text{USD }27,5/\text{MWh} \div 6,95\% = \mathbf{USD\text{ }486,6\text{ MM}}$, con un rango de USD 432,6 a 540,7 MM según factor de planta de 40\% a 50\%). Con $S_0 = \text{USD }486,6\text{ MM}$ y $K = I_0 = \text{USD }400\text{ MM}$, el algoritmo LSMC arroja un valor de \textbf{USD 186,9 MM} (+ARS 105,74 por acción / +USD 0,07).
    \item \textbf{Precio Objetivo Integrado (Base + Opción Real):} \textbf{ARS 1.374,22} por acción (USD 0,84, +39,9\% de retorno esperado).
\end{itemize}

\subsection{Prueba de Sensibilidad ante Repunte del Riesgo Soberano a 2.400 pb}
En un escenario de tensión macroeconómica donde el EMBI+ repunta a 2.400 pb, el WACC asciende a 13,21\% (+615 pb) y el precio objetivo se retrae a ARS 237,00 (Figura \ref{fig:17}).

\begin{figure}[H]
\centering
\includegraphics[width=0.65\linewidth,height=4.6cm,keepaspectratio]{../03_Modelo_y_Codigo/figuras/figura_17.png}
\caption{\textbf{Sensibilidad al Riesgo Soberano}: Impacto sobre WACC y Precio Objetivo en escenario de EMBI+ a 2.400 pb.}
\label{fig:17}
\end{figure}

\subsection{Escenarios}
\begin{table}[H]
\centering
\caption{\textbf{Marco de Escenarios: Supuestos Operativos y Precios Objetivos Resultantes}}
\scriptsize
\begin{tabularx}{\textwidth}{>{\raggedright\arraybackslash}p{2.5cm}>{\raggedright\arraybackslash}X>{\raggedright\arraybackslash}X>{\raggedright\arraybackslash}X}
\toprule
\textbf{Variable} & \textbf{Pesimista (25\%)} & \textbf{Base (50\%)} & \textbf{Optimista (25\%)} \\
\midrule
Precio LME & Caída hacia costo marginal global ($\sim$USD 1.900/t). & Nivel medio de equilibrio OU (USD 2.450--2.814/t). & LME sostenido $>$\,USD 3.000/t por déficit de oferta. \\
Riesgo País (EMBI+) & Rebote soberano a $>800$ pb. & Estabilización a largo plazo ($\approx$441 pb). & Compresión continua por debajo de 350 pb. \\
WACC / g & 9,0\% / 1,5\% & 6,95\% / 2,0\% & 6,5\% / 2,2\% \\
\midrule
\textbf{Precio Objetivo (DCF)} & \textbf{ARS 781} (-20,5\%) & \textbf{ARS 1.268,47} (Caso Base) & \textbf{ARS 1.481} (+50,7\%) \\
\bottomrule
\end{tabularx}
\end{table}

\subsection{Resumen de Precios Intrínsecos y Rangos Metodológicos}
La comparación de los distintos enfoques de valuación frente a la cotización Spot de mercado (Figura \ref{fig:18}) expone la dispersión del precio objetivo:

\begin{figure}[H]
\centering
\includegraphics[width=0.65\linewidth,height=4.6cm,keepaspectratio]{../03_Modelo_y_Codigo/figuras/figura_18.png}
\caption{\textbf{Rango de Precios Intrínsecos}: Comparativa de estimaciones metodológicas vs. cotización Spot.}
\label{fig:18}
\end{figure}

\begin{table}[H]
\centering
\caption{\textbf{Resumen de Precios Intrínsecos según Metodología (ARS/acción)}}
\small
\begin{tabularx}{\textwidth}{Xccc}
\toprule
\textbf{Metodología de Valuación} & \textbf{Mínimo} & \textbf{Central} & \textbf{Máximo} \\
\midrule
DCF $\lambda=1,0$ (Sin Cobertura de Riesgo País) & -- & 776 & -- \\
DCF con Tasa de Reinversión Damodaran & -- & 921 & -- \\
\textbf{DCF Caso Base Oficial (Dumrauf, $\lambda=0,1617$)} & -- & \textbf{1.268} & -- \\
DCF Crecimiento Optimista ($g=2,5\%$) & -- & 1.371 & -- \\
Monte Carlo (Percentiles P5 / Mediana / P95) & 866 & 1.270 & 1.790 \\
\midrule
\textbf{Precio Objetivo Integrado (DCF + Opción Real Etapa V Eólica)} & -- & \textbf{1.374,22} & -- \\
Escenarios Ponderados (25/50/25) & -- & 1.184 & -- \\
Cotización de Mercado & -- & 982,50 & -- \\
\bottomrule
\end{tabularx}
\end{table}

% ═════════════════════════════════════════════════════════════════════════
% SECCIÓN 12: ANÁLISIS DE SENSIBILIDAD DE VALUACIÓN
% ═════════════════════════════════════════════
\section{Matriz de Sensibilidad}

\subsection{Sensibilidad Cruzada de WACC y Tasa de Crecimiento Terminal ($g$)}
La matriz bidimensional (Figura \ref{fig:19}) confirma que la recomendación COMPRAR se mantiene válida para cualquier tasa WACC inferior a 8,0\% con tasas de crecimiento terminal $g \ge 1,5\%$.

\begin{figure}[H]
\centering
\includegraphics[width=0.65\linewidth,height=4.6cm,keepaspectratio]{../03_Modelo_y_Codigo/figuras/figura_19.png}
\caption{\textbf{Matriz WACC vs. g}: Sensibilidad bidimensional del precio fundamental por acción en ARS.}
\label{fig:19}
\end{figure}

% ═════════════════════════════════════════════════════════════════════════
% SECCIÓN 13: SIMULACIÓN MONTE CARLO Y DCF INVERSO
% ═════════════════════════════════════════════
\section{Simulación de Monte Carlo y DCF Inverso}

\subsection{Distribución Empírica con Innovaciones t-Student \texorpdfstring{($\nu=4,2$)}{(nu=4,2)}}
La simulación de Monte Carlo con 20.000 iteraciones e innovaciones con colas pesadas t-Student ($\nu = 4,2$, Figura \ref{fig:20}) aplica el factor de corrección $\sqrt{\frac{\nu-2}{\nu}}$:
\begin{equation}
    WACC_{sim} = WACC + \sigma_{WACC} \sqrt{\frac{\nu-2}{\nu}} \, T_\nu, \quad g_{sim} = g + \sigma_g \sqrt{\frac{\nu-2}{\nu}} \, T_\nu
\end{equation}
Resulta en una mediana de ARS 1.270 y un intervalo de confianza P5--P95 de ARS 866 a ARS 1.790, con un 87,7\% de probabilidad empírica de retorno positivo frente a la cotización de mercado.

\subsection{Tasa de Crecimiento Implícita del DCF Inverso \texorpdfstring{($g^* = 0,69\%$)}{(g* = 0,69\%)}}
Despejando la tasa de crecimiento terminal implícita en la cotización de mercado (Figura \ref{fig:21}), el precio descuenta un crecimiento a perpetuidad de $g^* = 0,69\%$, inferior al 2,0\% del modelo de equilibrio ($g$).

\begin{figure}[H]
\centering
\begin{minipage}[t]{0.485\textwidth}
    \centering
    \includegraphics[width=\linewidth,height=4.6cm,keepaspectratio]{../03_Modelo_y_Codigo/figuras/figura_20.png}
    \caption{\textbf{Simulación Monte Carlo (20.000 iteraciones)}: Distribución empírica con innovaciones t-Student ($\nu=4,2$).}
    \label{fig:20}
\end{minipage}\hfill
\begin{minipage}[t]{0.485\textwidth}
    \centering
    \includegraphics[width=\linewidth,height=4.6cm,keepaspectratio]{../03_Modelo_y_Codigo/figuras/figura_21.png}
    \caption{\textbf{DCF Inverso}: Tasa g implícita descontada en el precio de mercado (0,69\%).}
    \label{fig:21}
\end{minipage}
\end{figure}

% ═════════════════════════════════════════════════════════════════════════
% SECCIÓN 14: VALUACIÓN RELATIVA SECTORIAL
% ═════════════════════════════════════════════
\section{Valuación Relativa Sectorial y Múltiplos Comparables}

\subsection{Comparativa Internacional, Rendimiento FCF y Convergencia de Múltiplos}
El análisis comparativo sectorial confirma la divergencia entre la valuación basada en cifras históricas y la capacidad de generación normalizada. Mientras el múltiplo histórico de Aluar (13,4x) refleja la compresión de resultados de 2025, el rendimiento proyectado del flujo libre de caja ($>10\%$) supera ampliamente la mediana de los pares internacionales (3,8\%). A medida que se normaliza el EBITDA en 2026E--2027E, el múltiplo converge hacia 7,72x (Figura \ref{fig:22}):

\begin{figure}[H]
\centering
\includegraphics[width=0.65\linewidth,height=4.6cm,keepaspectratio]{../03_Modelo_y_Codigo/figuras/figura_22.png}
\caption{\textbf{Convergencia de Múltiplos}: Normalización del múltiplo EV/EBITDA hacia 7,72x.}
\label{fig:22}
\end{figure}

\begin{table}[H]
\centering
\caption{\textbf{Múltiplos Financieros y Rendimientos del Sector Aluminio (julio-2026)}}
\scriptsize
\begin{tabular}{>{\raggedright\arraybackslash}p{2.9cm} *4{>{\centering\arraybackslash}p{1.75cm}} >{\centering\arraybackslash}p{2.6cm}}
\toprule
\textbf{Compañía / Ticker} & \textbf{EV/EBITDA UDM} & \textbf{Crec. Ventas UDM} & \textbf{Margen EBITDA UDM} & \textbf{ROIC UDM} & \textbf{Rendimiento FCF 2026e/27e} \\
\midrule
ALUAR (ALUA.BA) & 13,4x & 8,0\% & 14,9\% & 4,3\% & -4,1\% / +12,8\% \\
Alcoa Corp. (AA) & 8,5x & 6,5\% & 13,1\% & 5,5\% & 4,8\% \\
Chalco (2600.HK) & 5,2x & 9,2\% & 15,5\% & 8,0\% & 5,2\% \\
Kaiser Aluminum (KALU) & 7,2x & 4,1\% & 9,8\% & 3,5\% & 3,0\% \\
Norsk Hydro (NHY.OL) & 7,7x & 5,8\% & 11,9\% & 6,2\% & 4,1\% \\
Constellium (CSTM) & 5,4x & 3,2\% & 10,1\% & 5,0\% & 3,5\% \\
Rusal & 5,1x & -2,0\% & 8,4\% & 2,1\% & 1,5\% \\
\midrule
\textbf{Mediana de Pares} & \textbf{6,3x} & \textbf{5,0\%} & \textbf{11,0\%} & \textbf{5,3\%} & \textbf{3,8\%} \\
\bottomrule
\end{tabular}
\end{table}

% ═════════════════════════════════════════════════════════════════════════
% SECCIÓN 15: GESTIÓN CUANTITATIVA DE RIESGO Y CARTERA
% ═════════════════════════════════════════════
\section{Gestión de Riesgos de Mercado y Asignación Óptima de Cartera}

\subsection{Criterio de Asignación de Kelly y Matriz de Riesgos Operativos}
El criterio de crecimiento de Kelly (Alexander, Sharpe \& Bailey, 2003) determina una fracción teórica óptima de:
\begin{equation}
    f^* = \frac{\mu - R_f}{\sigma^2} = \frac{24,63\% - 4,70\%}{(52,06\%)^2} = \mathbf{73,5\%} \text{ (Full-Kelly)} \implies Half-Kelly = \mathbf{36,8\%}
\end{equation}
Para la gestión prudencial de carteras se adopta un límite de exposición del 20,0\% (Figura \ref{fig:23}), acotado por el Conditional Value-at-Risk mensual al 95\% ($-32,80\%$). La categorización de los riesgos operativos y regulatorios se sintetiza en la matriz de probabilidad e impacto (Figura \ref{fig:28}).

\begin{figure}[H]
\centering
\begin{minipage}[t]{0.485\textwidth}
    \centering
    \includegraphics[width=\linewidth,height=4.6cm,keepaspectratio]{../03_Modelo_y_Codigo/figuras/figura_23.png}
    \caption{\textbf{Criterio de Asignación de Kelly}: Full-Kelly (73,5\%), Half-Kelly (36,8\%) y Límite de Cartera (20,0\%).}
    \label{fig:23}
\end{minipage}\hfill
\begin{minipage}[t]{0.485\textwidth}
    \centering
    \includegraphics[width=\linewidth,height=4.6cm,keepaspectratio]{../03_Modelo_y_Codigo/figuras/figura_28.png}
    \caption{\textbf{Matriz de Riesgos Operativos}: Clasificación por Probabilidad e Impacto de factores críticos.}
    \label{fig:28}
\end{minipage}
\end{figure}

\subsection{Comportamiento Histórico en Períodos de Estrés y Modelado EVT-GPD}
Ante la presencia de curtosis pronunciada (7,82) en los retornos en USD que invalida los supuestos gaussianos en eventos extremos, se calibra un modelo Pareto Generalizado (GPD, $\hat{\xi} = 0,214 > 0$, umbral $u$ al percentil 90\%), capturando pérdidas incondicionales de cola pesada Fréchet (Figura \ref{fig:25}):
\begin{equation}
    G_{\xi, \beta}(y) = 1 - \left(1 + \frac{\xi y}{\beta}\right)^{-1/\xi}
\end{equation}
Esta estimación cuantifica el riesgo de cola incondicional a 10 años ($\text{VaR}_{99\%} = -8,20\%$, $\text{ES}_{99\%} = -10,35\%$, Figura \ref{fig:25}). Por su parte, la simulación estocástica continua bajo procesos de difusión (Figura \ref{fig:27}) convalida la densidad de probabilidad del valor patrimonial con 10.000 trayectorias. El potencial de apreciación (+29,1\%) provee un margen de seguridad de más de 2,4 veces el déficit esperado centenario.

\begin{figure}[H]
\centering
\begin{minipage}[t]{0.485\textwidth}
    \centering
    \includegraphics[width=\linewidth,height=4.6cm,keepaspectratio]{../03_Modelo_y_Codigo/figuras/figura_25.png}
    \caption{\textbf{Valores Extremos (EVT-GPD)}: Estimación de colas pesadas con VaR 99\% (8,20\%) y ES 99\% (10,35\%).}
    \label{fig:25}
\end{minipage}\hfill
\begin{minipage}[t]{0.485\textwidth}
    \centering
    \includegraphics[width=\linewidth,height=4.6cm,keepaspectratio]{../03_Modelo_y_Codigo/figuras/figura_27.png}
    \caption{\textbf{DCF Estocástico Continuo}: Densidad de probabilidad del valor patrimonial con 10.000 trayectorias.}
    \label{fig:27}
\end{minipage}
\end{figure}

\begin{table}[H]
\centering
\caption{\textbf{VaR y CVaR Diario por Metodología (retornos en USD, horizonte 1 día)}}
\small
\begin{tabularx}{\textwidth}{Xccc}
\toprule
\textbf{Metodología de Riesgo} & \textbf{VaR 90\%} & \textbf{VaR 95\%} & \textbf{VaR 99\%} \\
\midrule
Histórico & -3,47\% & -4,73\% & -8,20\% \\
Normal (paramétrico) & -4,12\% & -5,31\% & -7,54\% \\
t-Student ($\nu=4,46$) & -3,61\% & -4,99\% & -8,57\% \\
Cornish-Fisher & -2,70\% & -5,22\% & -13,41\% \\
\textbf{EVT-GPD (Valores Extremos / Pareto)} & \textbf{-3,47\%} & \textbf{-4,73\%} & \textbf{-8,20\%} \\
CVaR / Expected Shortfall (Normal) & -5,67\% & -6,68\% & -8,65\% \\
\textbf{CVaR EVT-GPD (Pérdida Esperada Extrema)} & \textbf{-5,80\%} & \textbf{-7,15\%} & \textbf{-10,35\%} \\
\bottomrule
\end{tabularx}
\end{table}

\begin{table}[H]
\centering
\caption{\textbf{Desempeño Histórico de ALUA.BA (en USD) durante Episodios de Crisis}}
\scriptsize
\begin{tabularx}{\textwidth}{>{\raggedright\arraybackslash}p{3.6cm}ccc>{\raggedright\arraybackslash}X}
\toprule
\textbf{Episodio de Estrés} & \textbf{Retorno Acum.} & \textbf{Vol. Anualizada} & \textbf{Duración} & \textbf{Comportamiento y Dinámica} \\
\midrule
Crisis Covid-19 (Marzo 2020) & -44,73\% & 112,54\% & 50 ruedas & Shock de liquidez y contracción global. \\
Tensión Cambiaria (Julio 2022) & +12,44\% & 45,91\% & 103 ruedas & Cobertura eficiente en activos dolarizados. \\
Incertidumbre Electoral 2023 & +18,28\% & 105,36\% & 59 ruedas & Repreciación positiva de capacidad exportadora. \\
Corrección LME (2022--2023) & +13,61\% & 44,40\% & 200 ruedas & Resiliencia operativa por costos en primer cuartil. \\
\bottomrule
\end{tabularx}
\end{table}

% ═════════════════════════════════════════════════════════════════════════
% SECCIÓN 16: MODELIZACIÓN ESTOCÁSTICA Y ECONOMETRÍA
% ═════════════════════════════════════════════
\section{Modelos Estocásticos y Validación Econométrica}

\subsection{Procesos de Difusión del Commodity y la Tasa de Descuento}
Los commodities industriales no siguen un paseo aleatorio sin límite (MBG), sino que revierten a la media hacia el costo marginal de la industria (Figura \ref{fig:26}), modelado mediante un proceso Ornstein-Uhlenbeck:
\begin{equation}
    dS_t = \kappa_{\text{LME}} (\theta_{\text{LME}} - S_t) dt + \sigma_{\text{LME}} S_t dW_t, \quad \kappa = 0,384, \; \theta = \text{USD 2.814/t}
\end{equation}
Para el riesgo soberano (EMBI+), se calibra un proceso Cox-Ingersoll-Ross (CIR, Figura \ref{fig:29}) anclado al régimen macroeconómico de convergencia del DCF ($\theta = 441$ pb, $4,41\%$). Bajo el marco de cambio de régimen de Markov (Hamilton 1989, Gray 1996), la muestra histórica de 10 años ($N=2.594$ ruedas diarias) exhibe una media incondicional ergódica de $1.291$ pb debido a los episodios de estrés previos, pero la expectativa condicional en el régimen estabilizado actual converge a $\theta_1 = 441$ pb con $\kappa = 0,497$ y $\sigma = 18,70$ pb. El modelo satisface la condición de Feller ($2\kappa\theta > \sigma^2$), impidiendo tasas negativas y superando al AR(1) lineal por AIC ($\Delta AIC = -1.431,9$):
\begin{equation}
    dr_t = \kappa (\theta - r_t) dt + \sigma \sqrt{r_t} dW_t, \quad 2(0,497)(441) = 438,4 > 349,8 = \sigma^2
\end{equation}
Bajo el modelo de intensidad de crédito de Duffie-Singleton (1999) ($s_t = \lambda_t (1-R)$ con recupero esperado $R=35\%$), este spread equivale a una tasa de riesgo de default anual ($\lambda_t$) de $6,8\%$, frente al $19,9\%$ implícito en el promedio histórico de crisis.

\begin{figure}[H]
\centering
\begin{minipage}[t]{0.485\textwidth}
    \centering
    \includegraphics[width=\linewidth,height=4.6cm,keepaspectratio]{../03_Modelo_y_Codigo/figuras/figura_26.png}
    \caption{\textbf{Simulación LME (OU vs. MBG)}: Reversión a la media Ornstein-Uhlenbeck ($\theta=$ USD 2.814/t).}
    \label{fig:26}
\end{minipage}\hfill
\begin{minipage}[t]{0.485\textwidth}
    \centering
    \includegraphics[width=\linewidth,height=4.6cm,keepaspectratio]{../03_Modelo_y_Codigo/figuras/figura_29.png}
    \caption{\textbf{Proceso CIR Soberano}: Simulación condicional anclada a 441 pb con no-negatividad estricta.}
    \label{fig:29}
\end{minipage}
\end{figure}

\begin{table}[H]
\centering
\caption{\textbf{Bondad de Ajuste para el Modelado del Riesgo Soberano} ($N=2.594$ observaciones diarias, EMBI+ anclado a 441 pb)}
\small
\begin{tabularx}{\textwidth}{Xcc}
\toprule
\textbf{Modelo Estocástico} & \textbf{Log-Likelihood} & \textbf{AIC} \\
\midrule
AR(1) Lineal Base & -13.859,16 & 27.724,31 \\
\textbf{CIR (Difusión con Raíz Cuadrada)} & \textbf{-13.143,21} & \textbf{26.292,42} \\
\midrule
\multicolumn{3}{l}{\footnotesize $\Delta AIC = -1.431,89$ a favor del CIR (evidencia estadística concluyente).} \\
\bottomrule
\end{tabularx}
\end{table}

\subsection{Beta Dinámico (Filtro de Kalman, GARCH y Saltos)}
\label{subsec:kalman-garch}
El test ARCH-LM de rezago 5 ($LM = 291,8, p \approx 5,8\times10^{-61}$) confirma la presencia de heterocedasticidad condicional. Se estiman dos modelos dinámicos complementarios: el Filtro de Kalman (Figura \ref{fig:24}), que muestra la convergencia de $\beta_t$ a 0,851 en el régimen reciente, y el modelo GARCH(1,1) (Figura \ref{fig:30}):
\begin{equation}
    \sigma_t^2 = \omega + \alpha \epsilon_{t-1}^2 + \beta \sigma_{t-1}^2, \quad \hat{\alpha} + \hat{\beta} = 0,943, \quad y_t = \beta_t x_t + v_t
\end{equation}
Por el teorema de difusión límite de Nelson (1990), este GARCH(1,1) representa la aproximación discreta de un proceso de volatilidad estocástica continua de tipo Heston (1993). La descomposición de variación bipotencia (Barndorff-Nielsen \& Shephard 2004) bajo difusión con saltos de Merton (1976) muestra que el $28,4\%$ de la varianza histórica total proviene de shocks discretos devaluatorios, mientras que el Beta estructural de Hamada de 0,888 mantiene la calibración de ciclo completo para el WACC.

\begin{figure}[H]
\centering
\begin{minipage}[t]{0.485\textwidth}
    \centering
    \includegraphics[width=\linewidth,height=4.6cm,keepaspectratio]{../03_Modelo_y_Codigo/figuras/figura_24.png}
    \caption{\textbf{Beta Dinámico por Filtro de Kalman}: Convergencia a 0,851 en el régimen macroeconómico reciente.}
    \label{fig:24}
\end{minipage}\hfill
\begin{minipage}[t]{0.485\textwidth}
    \centering
    \includegraphics[width=\linewidth,height=4.6cm,keepaspectratio]{../03_Modelo_y_Codigo/figuras/figura_30.png}
    \caption{\textbf{Beta Condicional GARCH(1,1)}: Variación temporal de la volatilidad y riesgo sistemático.}
    \label{fig:30}
\end{minipage}
\end{figure}

\subsection{Selección de Cópulas y Sensibilidad de Sobol}
Se calibraron por máxima verosimilitud tres estructuras de cópulas sobre retornos diarios ALUA-USD vs. Merval-USD ($n=2.328$, dado que la serie EMBI+ local dispone de observaciones anuales): Gaussiana, Clayton $C_\theta(u,v)=(u^{-\theta}+v^{-\theta}-1)^{-1/\theta}$ y $t$-Student, comparándolas mediante el Criterio de Información de Akaike (AIC). La Cópula de Clayton ($\theta = 1,09$, $\lambda_L = 2^{-1/\theta} = 0,53$, Figura \ref{fig:31}) captura dependencia asimétrica en la cola inferior, pero es superada por la $t$-Student ($\rho = 0,623$, $\nu = 6,84$, $\lambda_L = \lambda_U = 0,215$), que registra mayor bondad de ajuste relativo por AIC frente a Clayton ($\Delta AIC = +207,9$) y frente a la Gaussiana ($\Delta AIC = +161,5$). Esto muestra que la comovilidad entre el activo y el índice bursátil doméstico exhibe colas pesadas simétricas, donde ambos extremos (movimientos conjuntos de gran magnitud al alza o a la baja) tienen probabilidades equiparables, descartando una asimetría bajista predominante.

\begin{figure}[H]
\centering
\begin{minipage}[t]{0.485\textwidth}
    \centering
    \includegraphics[width=\linewidth,height=4.6cm,keepaspectratio]{../03_Modelo_y_Codigo/figuras/figura_31.png}
    \caption{\textbf{Ajuste de Cópula de Clayton}: Ilustra la dependencia de cola inferior ($\lambda_L = 0,53$). El ajuste por AIC favorece finalmente a la $t$-Student (cola simétrica, ver texto).}
    \label{fig:31}
\end{minipage}\hfill
\begin{minipage}[t]{0.485\textwidth}
    \begin{tcolorbox}[colback=white,colframe=gsnavy,arc=2pt,boxrule=0.8pt,top=3pt,bottom=3pt,left=3pt,right=3pt]
        \centering
        {\small\bfseries\color{gsnavy} Descomposición de Varianza (Sobol)}
        \vspace{2pt}
        \scriptsize
        \begin{tabularx}{\linewidth}{Xrr}
            \toprule
            \textbf{Variable} & \textbf{Índice $S_i$} & \textbf{Total $S_{Ti}$} \\
            \midrule
            Shock de margen EBITDA & 0,395 & 0,427 \\
            $g$ (crecimiento perpetuo) & 0,296 & 0,377 \\
            WACC & 0,187 & 0,263 \\
            \bottomrule
        \end{tabularx}
    \end{tcolorbox}
    \vspace{2pt}
    {\scriptsize Sobre las mismas tres fuentes de incertidumbre del Monte Carlo (Figura \ref{fig:20}). Los índices de Sobol muestran que el supuesto de margen EBITDA de largo plazo explica más varianza del precio que $g$ o el WACC.}
\end{minipage}
\end{figure}

\subsection{Diagnóstico Econométrico y Convalidación del Modelo de Valuación}
La validez del modelo fundamental descansa en la verificación econométrica de sus supuestos. La ventana muestral de 10 años (2016--2026, $N=2.344$ retornos diarios) fue seleccionada para capturar un ciclo completo de materias primas y riesgo soberano, evitando la distorsión estructural de regímenes previos (Convertibilidad y defaults 2001--2005) donde la capacidad instalada y la matriz energética de Aluar eran sustancialmente distintas. El test de estabilidad de Quandt-Andrews sobre la regresión de mercado ($\sup F = 11,27$, $p \approx 0,00$ por bootstrap de residuos, recorte 15\,\%--85\,\%) sí detecta un quiebre paramétrico significativo, fechado el 19 de marzo de 2020 (semana del colapso global por COVID-19): el Beta salta de 1,208 en el tramo previo a 0,665 en el posterior. Se conserva de todos modos la regresión de 10 años sin segmentar porque (i) es la ventana estándar de referencia para el Beta OLS en la práctica de valuación, (ii) el quiebre corresponde a un evento de mercado extremo y transitorio, no a un cambio estructural del negocio de Aluar, y (iii) el propio proceso de Blume y el Filtro de Kalman de la Sección \ref{subsec:kalman-garch} ya corrigen el Beta hacia su nivel reciente, absorbiendo en la práctica el efecto del quiebre.

\begin{table}[H]
\centering
\caption{\textbf{Resumen de Pruebas Estadísticas y Convalidación del Modelo Base}}
\scriptsize
\begin{tabularx}{\textwidth}{>{\raggedright\arraybackslash}p{3.2cm}>{\raggedright\arraybackslash}p{2.8cm}cc>{\raggedright\arraybackslash}X}
\toprule
\textbf{Prueba Econométrica} & \textbf{Variable / Proceso} & \textbf{Estadístico} & \textbf{P-Valor} & \textbf{Conclusión y Respaldo a la Valuación} \\
\midrule
Dickey-Fuller Aumentado (ADF) & Retornos Diarios ALUA (USD) & $t = -10,44$ & $p \approx 0$ & Serie estacionaria $I(0)$, descarta derivas espurias. \\
Phillips-Perron (PP) & Retornos Diarios ALUA (USD) & $Z_t = -47,37$ & $p \approx 0$ & Estacionario bajo heterocedasticidad y autocorrelación. \\
Jarque-Bera & Normalidad de Retornos & $JB = 4.391,7$ & $p < 0,0001$ & Rechazo de normalidad (Curtosis = 7,82), valida EVT. \\
Kolmogorov-Smirnov & t-Student ($\nu=4,46$) vs. Normal & $D = 0,018$ & $p = 0,42$ & Convalida colas pesadas en simulación Monte Carlo. \\
Cópula ($t$-Student vs. Clayton) & Dependencia ALUA--Merval & $\Delta AIC = +207,9$ & $\lambda_L=\lambda_U=0,215$ & $t$-Student preferida, descarta colapso unilateral. \\
ARCH-LM (Lag 5) & Heterocedasticidad Condicional & $LM = 291,8$ & $p \approx 0$ & Agrupamiento de volatilidad modelado por GARCH(1,1). \\
Filtro de Kalman & Beta Dinámico $\beta_t$ & $\beta_{\text{reciente}} = 0,851$ & --- & Convalida el Beta estructural Hamada ($\beta_L = 0,888$). \\
CIR vs. AR(1) Lineal & Difusión del Riesgo Soberano & $\Delta AIC = -1.431,9$ & $2\kappa\theta > \sigma^2$ & Impide tasas negativas en estrés y acota el WACC. \\
Quandt-Andrews & Estabilidad de Parámetros (Beta OLS) & $\sup F = 11,27$ & $p \approx 0,00$ & Detecta quiebre en 19-mar-2020 (COVID). $\beta$ 1,208 $\to$ 0,665, corregido por Blume/Kalman. \\
\bottomrule
\end{tabularx}
\end{table}

% ═════════════════════════════════════════════════════════════════════════
% SECCIÓN 17: CONCLUSIÓN
% ═════════════════════════════════════════════
\section{Dictamen del Modelo de Valuación y Conclusiones}

El dictamen técnico determina \textbf{COMPRAR} sobre los siguientes fundamentos:
\begin{itemize}
    \item \textbf{Precio Objetivo Base (DCF Dumrauf):} \textbf{ARS 1.268,47} por acción (USD 0,80, retorno esperado de +29,1\%).
    \item \textbf{Precio Objetivo Integrado (Base + Opción Real Etapa V Eólica):} \textbf{ARS 1.374,22} por acción (USD 0,84, retorno esperado de +39,9\%). El activo subyacente $S_0$ surge del valor presente de los ahorros de energía a costo WACC con factor de planta del 45\%.
    \item \textbf{Rango Probabilístico Monte Carlo (P5--P95):} ARS 866 a ARS 1.790 por acción, con 87,7\% de probabilidad de retorno positivo.
    \item \textbf{Margen de Seguridad:} La cotización actual descuenta una tasa terminal implícita de $g^* = 0,69\%$, con un margen de ARS 253,50 por acción frente al valor base.
\end{itemize}

\clearpage

% ═════════════════════════════════════════════════════════════════════════
% APÉNDICE: ESTADOS FINANCIEROS AUDITADOS Y PROYECCIONES
% ═════════════════════════════════════════════
\section*{Anexo de Estados Financieros Auditados y Proyecciones FCFF (USD MM)}
\addcontentsline{toc}{section}{Anexo de Estados Financieros Auditados y Proyecciones FCFF}

{\small Cifras históricas anuales auditadas por Price Waterhouse \& Co. S.R.L. convertidas al CCL de cierre de cada ejercicio. Las proyecciones aplican la alícuota estatutaria del 35\% sobre el resultado operativo (Ley 20.628).}

\begin{table}[H]
\centering
\caption{\textbf{Conciliación de la Tasa Efectiva del Impuesto a las Ganancias 2025}}
\small
\begin{tabularx}{\textwidth}{Xcc}
\toprule
\textbf{Componente Tributario} & \textbf{Monto (USD MM)} & \textbf{\% sobre EBT} \\
\midrule
Resultado Antes de Impuestos (EBT) & 58,5 & 100,0\% \\
Impuesto Teórico a Tasa Estatutaria (35\%) & (20,5) & 35,0\% \\
Efecto del Ajuste por Inflación Impositivo (NIC 29) & (23,1) & 39,5\% \\
Diferencias de Cambio no Deducibles y Otros & (5,7) & 9,8\% \\
\midrule
\textbf{Impuesto a las Ganancias Devengado en Estados Contables} & \textbf{(49,3)} & \textbf{84,3\%} \\
\bottomrule
\end{tabularx}
\end{table}

\begin{table}[H]
\centering
\caption{\textbf{Estado de Resultados Consolidado Auditado (2020--2025, USD MM)}}
\scriptsize
\begin{tabularx}{\textwidth}{Xrrrrrr}
\toprule
\textbf{Línea Contable} & \textbf{2020} & \textbf{2021} & \textbf{2022} & \textbf{2023} & \textbf{2024} & \textbf{2025} \\
\midrule
Ventas Netas (Revenue) & 823,5 & 513,5 & 654,7 & 647,8 & 915,9 & 1.092,6 \\
Costo de Ventas & (711,2) & (414,4) & (448,3) & (536,0) & (744,0) & (930,5) \\
\textbf{Ganancia Bruta} & \textbf{112,3} & \textbf{99,1} & \textbf{206,4} & \textbf{111,8} & \textbf{171,9} & \textbf{162,1} \\
Gastos de Administración y Comercialización & (63,8) & (37,8) & (44,1) & (50,1) & (70,1) & (100,8) \\
Otros Ingresos / Operativos Netos & (0,4) & 0,0 & (0,2) & 0,4 & 50,8 & 31,0 \\
\textbf{EBITDA} & \textbf{116,5} & \textbf{110,2} & \textbf{208,8} & \textbf{103,5} & \textbf{204,7} & \textbf{163,3} \\
Depreciación y Amortización (D\&A) & (68,4) & (48,9) & (46,7) & (41,5) & (52,1) & (71,0) \\
\textbf{EBIT (Resultado Operativo)} & \textbf{48,1} & \textbf{61,3} & \textbf{162,1} & \textbf{62,0} & \textbf{152,6} & \textbf{92,3} \\
Resultado Financiero y Asociadas & (84,5) & 14,7 & 36,2 & 88,4 & 3,9 & (33,8) \\
\textbf{Ganancia Antes de Impuestos (EBT)} & \textbf{(36,4)} & \textbf{76,0} & \textbf{198,3} & \textbf{150,3} & \textbf{156,5} & \textbf{58,5} \\
Impuesto a las Ganancias & (7,5) & (47,9) & (82,8) & (11,3) & (66,7) & (49,3) \\
\textbf{Resultado Neto del Ejercicio} & \textbf{(43,9)} & \textbf{28,1} & \textbf{115,5} & \textbf{139,0} & \textbf{89,7} & \textbf{9,2} \\
\textbf{NOPAT (Resultado Operativo Neto de Impuestos)} & \textbf{31,3} & \textbf{39,9} & \textbf{105,4} & \textbf{40,3} & \textbf{99,2} & \textbf{60,0} \\
\bottomrule
\end{tabularx}
\end{table}

\begin{table}[H]
\centering
\caption{\textbf{Estado de Situación Patrimonial / Balance Consolidado Auditado (2020--2025, USD MM)}}
\scriptsize
\begin{tabularx}{\textwidth}{Xrrrrrr}
\toprule
\textbf{Rubro Patrimonial} & \textbf{2020} & \textbf{2021} & \textbf{2022} & \textbf{2023} & \textbf{2024} & \textbf{2025} \\
\midrule
Activo Corriente & 454,6 & 318,1 & 478,9 & 490,1 & 902,9 & 1.028,1 \\
Activo No Corriente (inc. Bienes de Uso) & 597,4 & 379,9 & 353,6 & 422,3 & 556,6 & 940,9 \\
\textbf{Total Activo} & \textbf{1.052,0} & \textbf{698,0} & \textbf{832,6} & \textbf{912,4} & \textbf{1.459,5} & \textbf{1.969,0} \\
Pasivo Corriente & 240,1 & 84,2 & 163,8 & 167,8 & 183,2 & 420,5 \\
Pasivo No Corriente & 348,7 & 263,9 & 212,2 & 201,9 & 435,9 & 439,2 \\
\textbf{Total Pasivo} & \textbf{588,8} & \textbf{348,1} & \textbf{376,0} & \textbf{369,7} & \textbf{619,1} & \textbf{859,7} \\
Deuda Financiera Total & 375,2 & 181,4 & 134,0 & 228,9 & 408,5 & 594,9 \\
Caja, Bancos e Inversiones & 23,9 & 44,0 & 76,7 & 59,0 & 197,3 & 69,2 \\
\textbf{Deuda Financiera Neta} & \textbf{351,3} & \textbf{137,4} & \textbf{57,3} & \textbf{169,9} & \textbf{211,2} & \textbf{525,8} \\
\textbf{Patrimonio Neto} & \textbf{463,2} & \textbf{349,9} & \textbf{456,5} & \textbf{542,7} & \textbf{840,4} & \textbf{1.109,3} \\
\textbf{Capital Invertido (NOA)} & \textbf{838,4} & \textbf{531,3} & \textbf{590,5} & \textbf{771,6} & \textbf{1.248,9} & \textbf{1.704,3} \\
\bottomrule
\end{tabularx}
\end{table}

\begin{table}[H]
\centering
\caption{\textbf{Estado de Flujos de Efectivo Consolidado Auditado (2020--2025, USD MM)}}
\scriptsize
\begin{tabularx}{\textwidth}{Xrrrrrr}
\toprule
\textbf{Flujo de Caja} & \textbf{2020} & \textbf{2021} & \textbf{2022} & \textbf{2023} & \textbf{2024} & \textbf{2025} \\
\midrule
Flujo de Caja Operativo (FCO) & 245,7 & 102,4 & 89,2 & 113,6 & 79,2 & 30,7 \\
Inversiones en Bienes de Uso (CAPEX) & (103,0) & (9,4) & (11,1) & (65,9) & (25,6) & (243,9) \\
Flujo de Caja de Inversión (FCI) & (102,1) & (9,4) & (11,1) & (65,9) & (25,6) & (244,0) \\
Flujo de Caja de Financiación (FCF) & (197,8) & (45,1) & (34,4) & (70,8) & 91,1 & 20,3 \\
\textbf{Variación Neta de Efectivo} & \textbf{(54,2)} & \textbf{47,8} & \textbf{43,8} & \textbf{(23,1)} & \textbf{144,7} & \textbf{(193,0)} \\
\bottomrule
\end{tabularx}
\end{table}

\begin{table}[H]
\centering
\caption{\textbf{Proyecciones del Free Cash Flow to the Firm (FCFF, 2026e--2030e, USD MM)}}
\scriptsize
\begin{tabularx}{\textwidth}{Xrrrrr}
\toprule
\textbf{Cuenta / Ejercicio} & \textbf{2026E} & \textbf{2027E} & \textbf{2028E} & \textbf{2029E} & \textbf{2030E} \\
\midrule
Ventas Netas (Revenue) & 1.591,4 & 1.543,0 & 1.494,6 & 1.446,3 & 1.397,9 \\
EBITDA & 391,2 & 369,3 & 347,4 & 325,5 & 303,6 \\
(-) Depreciación y Amortización & (108,4) & (105,1) & (101,8) & (98,5) & (95,2) \\
\textbf{EBIT} & \textbf{282,8} & \textbf{264,2} & \textbf{245,6} & \textbf{227,0} & \textbf{208,4} \\
(-) Impuestos Operativos (35\% Estatutaria) & (99,0) & (92,5) & (86,0) & (79,4) & (72,9) \\
\textbf{NOPAT} & \textbf{183,8} & \textbf{171,7} & \textbf{159,6} & \textbf{147,6} & \textbf{135,5} \\
(+) D\&A (Reversión) & 108,4 & 105,1 & 101,8 & 98,5 & 95,2 \\
(-) CAPEX de Mantenimiento & (103,2) & (100,0) & (96,9) & (93,8) & (90,6) \\
(-) Variación Capital de Trabajo ($\Delta$NWC) & (246,5) & 23,9 & 23,9 & 23,9 & 23,9 \\
\midrule
\textbf{FCFF Explícito} & \textbf{-57,5} & \textbf{200,7} & \textbf{188,4} & \textbf{176,2} & \textbf{164,0} \\
\bottomrule
\end{tabularx}
\end{table}

\noindent
{\small El FCFF Explícito 2030E de USD 164,0 MM (con CAPEX de mantenimiento de USD 90,6 MM y liberación de NWC de USD 23,9 MM) se reconcilia con el FCFF Terminal Normalizado ($g = 2,0\%$) de USD 135,5 MM, que asume estado estacionario donde CAPEX = D\&A y $\Delta$NWC = \$0, coincidiendo con el NOPAT de 2030E según la metodología de Dumrauf Cap. 14.\par}

% ═════════════════════════════════════════════════════════════════════════
% ANEXO DE FUENTES Y METODOLOGÍA
% ═════════════════════════════════════════════
\section*{Anexo de Fuentes y Metodología}
\addcontentsline{toc}{section}{Anexo de Fuentes y Metodología}

\noindent
\textbf{\color{gsnavy} Fuentes de Información y Fecha de Corte}\\[2pt]
{\small Las estimaciones y proyecciones se basan únicamente en información pública: Estados Financieros Consolidados Auditados de Aluar (dictamen de Price Waterhouse \& Co. S.R.L., ejercicios cerrados al 30 de junio 2020--2025), Memorias Anuales, hechos relevantes publicados en CNV y BYMA, estadísticas del London Metal Exchange (LME), datos del BCRA y series de la Reserva Federal (FRED). Las variables de mercado (LME spot, EMBI+ Argentina, CCL implícito vía GGAL) corresponden al cierre del 23 de julio de 2026. No reflejan movimientos posteriores a esa fecha.\par}

\vspace{0.2cm}
\noindent
\textbf{\color{gsnavy} Metodología Cuantitativa y Supuestos Clave}\\[2pt]
{\small El precio objetivo surge de un modelo de Flujo de Fondos Descontados (DCF) bajo la formulación de Dumrauf (Cap. 14) con WACC ajustado por riesgo país ($\lambda=0,1617$), Beta Hamada, crecimiento terminal $g=2,0\%$, complementado con opción real de expansión de la Etapa V eólica (LSMC Longstaff-Schwartz), simulación Monte Carlo (20.000 iteraciones, innovaciones t-Student $\nu=4,2$), EVT-GPD, GARCH(1,1), Filtro de Kalman, y selección de cópulas por AIC. Horizonte de inversión de referencia: 12 meses. Índice de referencia: S\&P Merval (MERVAL). Escala de recomendación: COMPRAR (retorno esperado $>+15\%$), MANTENER ($-15\%$ a $+15\%$), VENDER ($<-15\%$).\par}

\vspace{0.2cm}
\noindent
\textbf{\color{gsnavy} Vigencia y Próxima Revisión}\\[2pt]
{\small Fecha de emisión: 23 de julio de 2026. Próxima revisión programada: 23 de octubre de 2026 (frecuencia trimestral) o ante evento material que modifique sustancialmente los supuestos de valuación (resultados trimestrales, cambios macroeconómicos significativos, anuncios corporativos relevantes).\par}

% ═════════════════════════════════════════════════════════════════════════
% REFERENCIAS BIBLIOGRÁFICAS
% ═════════════════════════════════════════════
\section*{Bibliografía y Referencias}
\addcontentsline{toc}{section}{Bibliografía y Referencias}

\begin{multicols}{2}
\begin{enumerate}
    \tiny
    \item Alexander, G.J., Sharpe, W.F., Bailey, J.V. (2003). \textit{Fundamentos de Inversiones: Teoría y Práctica}, 3ª Ed., Pearson Educación, México.
    \item Aluar Aluminio Argentino S.A.I.C. (2020–2025). \textit{Estados Financieros Consolidados Auditados y Memorias Anuales}, CNV \& BYMA.
    \item Barndorff-Nielsen, O.E., Shephard, N. (2004). ``Power and Bipower Variation with Stochastic Volatility and Jumps'', \textit{J. Financ. Econometrics}, 2(1), pp. 1--37.
    \item Black, F., Scholes, M. (1973). ``The Pricing of Options and Corporate Liabilities'', \textit{J. Polit. Econ.}, 81(3), pp. 637--654.
    \item Blume, M.E. (1975). ``Betas and Their Regression Tendencies'', \textit{J. Finance}, 30(3), pp. 785--795.
    \item Bollerslev, T. (1986). ``Generalized Autoregressive Conditional Heteroskedasticity'', \textit{J. Econometrics}, 31(3), pp. 307--327.
    \item Cox, J.C., Ingersoll, J.E., Ross, S.A. (1985). ``A Theory of the Term Structure of Interest Rates'', \textit{Econometrica}, 53(2), pp. 385--407.
    \item Damodaran, A. (2012). \textit{Investment Valuation}, 3rd Ed., John Wiley \& Sons.
    \item Duffie, D., Singleton, K.J. (1999). ``Modeling Term Structures of Defaultable Bonds'', \textit{Rev. Financ. Stud.}, 12(4), pp. 687--720.
    \item Dumrauf, G.L. (2013). \textit{Finanzas Corporativas: Un Enfoque Latinoamericano}, 3ª Ed., Alfaomega.
    \item Engle, R.F. (1982). ``Autoregressive Conditional Heteroscedasticity'', \textit{Econometrica}, 50(4), pp. 987--1007.
    \item Gray, S.F. (1996). ``Modeling the Conditional Distribution of Interest Rates as a Regime-Switching Process'', \textit{J. Financ. Econ.}, 42(1), pp. 27--62.
    \item Hamada, R.S. (1972). ``The Effect of the Firm's Capital Structure on Systematic Risk'', \textit{J. Finance}, 27(2), pp. 435--452.
    \item Hamilton, J.D. (1989). ``A New Approach to the Economic Analysis of Nonstationary Time Series'', \textit{Econometrica}, 57(2), pp. 357--384.
    \item Heston, S.L. (1993). ``A Closed-Form Solution for Options with Stochastic Volatility'', \textit{Rev. Financ. Stud.}, 6(2), pp. 327--343.
    \item Jaschke, S.R. (2001). ``The Cornish-Fisher-Expansion in Delta-Gamma-Normal Approximations'', \textit{J. Risk}, 4(4), pp. 33--52.
    \item Kelly, J.L. (1956). ``A New Interpretation of Information Rate'', \textit{Bell Syst. Tech. J.}, 35(4), pp. 917--926.
    \item Longstaff, F.A., Schwartz, E.S. (2001). ``Valuing American Options by Simulation: A Simple Least-Squares Approach'', \textit{Rev. Financ. Stud.}, 14(1), pp. 113--147.
    \item McNeil, A.J., Frey, R., Embrechts, P. (2015). \textit{Quantitative Risk Management}, Revised Ed., Princeton University Press.
    \item Merton, R.C. (1976). ``Option Pricing when Underlying Stock Returns Are Discontinuous'', \textit{J. Financ. Econ.}, 3(1-2), pp. 125--144.
    \item Nelson, D.B. (1990). ``ARCH Models as Diffusion Approximations'', \textit{J. Econometrics}, 45(1-2), pp. 7--38.
    \item Rockafellar, R.T., Uryasev, S. (2000). ``Optimization of Conditional Value-at-Risk'', \textit{J. Risk}, 2(3), pp. 21--42.
    \item Ross, S.A., Westerfield, R.W., Jaffe, J. (2019). \textit{Corporate Finance}, 12th Ed., McGraw-Hill.
    \item Sklar, A. (1959). ``Fonctions de Répartition à n Dimensions et Leurs Marges'', \textit{Publ. Inst. Statist. Univ. Paris}, 8, pp. 229--231.
    \item Sobol, I.M. (2001). ``Global Sensitivity Indices for Nonlinear Mathematical Models'', \textit{Math. Comput. Simulation}, 55(1-3), pp. 271--280.
\end{enumerate}
\end{multicols}

\end{document}
